\documentclass[10pt,a4paper]{article}

% ─── Packages ────────────────────────────────────────────────────────────────
\usepackage[a4paper, top=2cm, bottom=2cm, left=2.2cm, right=2.2cm]{geometry}
\usepackage{amsmath, amssymb, amsthm}
\usepackage{listings}
\usepackage{xcolor}
\usepackage{hyperref}
\usepackage{titlesec}
\usepackage{parskip}
\usepackage{tocloft}
\usepackage{fancyhdr}
\usepackage{mdframed}
\usepackage{enumitem}
\usepackage{booktabs}
\usepackage{multicol}
\usepackage{microtype}
\usepackage{array}

% ─── Colors ──────────────────────────────────────────────────────────────────
\definecolor{codebg}{RGB}{245,245,250}
\definecolor{codecomment}{RGB}{80,130,80}
\definecolor{codekw}{RGB}{0,80,180}
\definecolor{codestr}{RGB}{180,0,0}
\definecolor{accent}{RGB}{25,80,160}
\definecolor{theorembg}{RGB}{240,248,255}
\definecolor{notebg}{RGB}{255,250,235}
\definecolor{sectioncolor}{RGB}{25,80,160}

% ─── Listings (Python) ───────────────────────────────────────────────────────
\lstdefinestyle{pycode}{
  language=Python,
  backgroundcolor=\color{codebg},
  basicstyle=\ttfamily\small,
  keywordstyle=\color{codekw}\bfseries,
  stringstyle=\color{codestr},
  commentstyle=\color{codecomment}\itshape,
  numberstyle=\tiny\color{gray},
  numbers=left,
  numbersep=8pt,
  stepnumber=1,
  showstringspaces=false,
  breaklines=true,
  breakatwhitespace=false,
  frame=single,
  framesep=4pt,
  rulecolor=\color{gray!40},
  tabsize=4,
  captionpos=b,
  xleftmargin=15pt,
  xrightmargin=5pt,
  aboveskip=8pt,
  belowskip=8pt,
  morekeywords={True, False, None, yield, lambda, assert, with, as, pass},
}
\lstset{style=pycode}

% ─── Theorem Environments ────────────────────────────────────────────────────
\newtheoremstyle{cpstyle}
  {8pt}{8pt}{\normalfont}{0pt}{\bfseries\color{accent}}{.}{.5em}{}

\theoremstyle{cpstyle}
\newtheorem{theorem}{Theorem}[section]
\newtheorem{lemma}[theorem]{Lemma}
\newtheorem{corollary}[theorem]{Corollary}
\newtheorem{definition}[theorem]{Definition}
\newtheorem{algorithm}[theorem]{Algorithm}
\newtheorem{observation}[theorem]{Observation}

\newenvironment{cpnote}{%
  \begin{mdframed}[backgroundcolor=notebg, linecolor=orange!70, linewidth=1pt,
    innerleftmargin=10pt, innerrightmargin=10pt, innertopmargin=6pt, innerbottommargin=6pt]
  \textbf{\color{orange!80!black}Note / Pitfall:}~}{%
  \end{mdframed}}

\newenvironment{cputility}{%
  \begin{mdframed}[backgroundcolor=theorembg, linecolor=accent!50, linewidth=1pt,
    innerleftmargin=10pt, innerrightmargin=10pt, innertopmargin=6pt, innerbottommargin=6pt]
  \textbf{\color{accent}CP Utility:}~}{%
  \end{mdframed}}

% ─── Section Formatting ──────────────────────────────────────────────────────
\newcommand{\sectionbox}[1]{%
  \colorbox{accent}{\parbox{\dimexpr\linewidth-2\fboxsep\relax}{\thesection\quad #1}}}

\titleformat{\section}[block]
  {\Large\bfseries\color{white}}
  {}
  {0pt}
  {\sectionbox}
  [\vspace{2pt}]

\titleformat{\subsection}
  {\large\bfseries\color{accent}}
  {\thesubsection\quad}{0pt}{}
  [\titlerule\vspace{2pt}]

\titleformat{\subsubsection}
  {\normalsize\bfseries\color{gray!70!black}}
  {\thesubsubsection\quad}{0pt}{}

% ─── Header / Footer ─────────────────────────────────────────────────────────
\pagestyle{fancy}
\fancyhf{}
\fancyhead[L]{\small\textcolor{accent}{\textbf{Winter Cup / ICPC Handbook}}}
\fancyhead[R]{\small\textcolor{gray}{\leftmark}}
\fancyfoot[C]{\small\thepage}
\renewcommand{\headrulewidth}{0.4pt}
\renewcommand{\footrulewidth}{0.2pt}

% ─── Hyperref Setup ──────────────────────────────────────────────────────────
\hypersetup{
  colorlinks=true,
  linkcolor=accent,
  urlcolor=blue!70!black,
  citecolor=green!60!black,
  pdftitle={Winter Cup / ICPC Competitive Programming Handbook},
  pdfauthor={CP Handbook},
  pdfsubject={Competitive Programming},
  bookmarksnumbered=true,
}

% ─── Misc Helpers ────────────────────────────────────────────────────────────
\newcommand{\bigO}[1]{\mathcal{O}\!\left(#1\right)}
\newcommand{\Z}{\mathbb{Z}}
\newcommand{\N}{\mathbb{N}}
\newcommand{\R}{\mathbb{R}}

% ─────────────────────────────────────────────────────────────────────────────
\begin{document}

% ─────────────────────────────────────────────────────────────────────────────
%  TITLE PAGE
% ─────────────────────────────────────────────────────────────────────────────
\begin{titlepage}
  \begin{center}
    \vspace*{2cm}
    {\Huge\bfseries\color{accent} Winter Cup}\\[0.5em]
    {\LARGE\bfseries Competitive Programming Handbook}\\[1em]
    {\large ICPC-Style Contest Reference}\\[3em]
    \hrule height 2pt
    \vspace{1em}
    {\large\itshape Theorems $\cdot$ Algorithms $\cdot$ Implementations $\cdot$ Proofs}\\[0.5em]
    \hrule height 2pt
    \vspace{3em}
    \begin{mdframed}[backgroundcolor=theorembg, linecolor=accent, linewidth=2pt,
      innerleftmargin=20pt, innerrightmargin=20pt, innertopmargin=16pt, innerbottommargin=16pt]
    \textbf{Contents:}
    \begin{multicols}{2}
    \begin{itemize}[noitemsep]
      \item[\S1] Number Theory
      \item[\S2] Combinatorics
      \item[\S3] Graph Theory
      \item[\S4] Dynamic Programming
      \item[\S5] Greedy Algorithms
      \item[\S6] Data Structures
      \item[\S7] String Algorithms
      \item[\S8] Geometry
      \item[\S9] Probability \& Math
      \item[\S10] CP Patterns
    \end{itemize}
    \end{multicols}
    \end{mdframed}
    \vfill
    {\small All implementations in Python 3 $\cdot$ Contest-ready $\cdot$ PDF printable}\\[0.5em]
    {\small\color{gray} Compiled for ICPC / ACM-style contests}
  \end{center}
\end{titlepage}

% ─────────────────────────────────────────────────────────────────────────────
%  TABLE OF CONTENTS
% ─────────────────────────────────────────────────────────────────────────────
\setcounter{tocdepth}{2}
\tableofcontents
\clearpage

% =============================================================================
\section{Number Theory}
% =============================================================================

% ─────────────────────────────────────────────────────────────────────────────
\subsection{Fermat's Little Theorem}

\begin{theorem}[Fermat's Little Theorem]
Let $p$ be a prime and $a$ an integer with $\gcd(a,p)=1$. Then
\[
  a^{p-1} \equiv 1 \pmod{p}.
\]
Equivalently, $a^p \equiv a \pmod{p}$ for every integer $a$.
\end{theorem}

\textbf{Intuition.} When you raise $a$ to the power $p-1$ modulo a prime $p$, the group structure of $(\Z/p\Z)^*$ forces the result to collapse to $1$.

\begin{cputility}
Fast computation of modular inverses: $a^{-1} \equiv a^{p-2} \pmod{p}$ when $p$ is prime.
\end{cputility}

\textbf{Where to use.}
\begin{itemize}[noitemsep]
  \item Computing $\binom{n}{k} \bmod p$ (modular inverse of factorials).
  \item Reducing huge exponents: $a^x \equiv a^{x \bmod (p-1)} \pmod{p}$.
  \item Hashing, cryptography problems.
\end{itemize}

\textbf{Typical CP scenarios.} Counting paths, number of ways, divide-and-conquer with modular arithmetic.

\textbf{Example Problem.} Compute $\binom{n}{k} \bmod p$ for large $n,k$ where $p$ is prime.

\subsubsection{Python Implementation}
\begin{lstlisting}[caption={Fermat's Little Theorem -- Modular Inverse}]
MOD = 10**9 + 7

def mod_inv(a, p=MOD):
    # Fermat: a^(p-2) mod p = a^(-1) mod p
    return pow(a, p - 2, p)

def comb_mod(n, k, p=MOD):
    if k < 0 or k > n:
        return 0
    num = 1
    den = 1
    for i in range(k):
        num = num * (n - i) % p
        den = den * (i + 1) % p
    return num * mod_inv(den, p) % p

# Example
print(comb_mod(10, 3))  # C(10,3) mod 10^9+7 = 120
\end{lstlisting}

\textbf{Complexity.} $\bigO{\log p}$ per inverse via fast exponentiation.

\begin{cpnote}
Only valid when $p$ is prime. For composite moduli use the Extended Euclidean Algorithm instead.
\end{cpnote}

% ─────────────────────────────────────────────────────────────────────────────
\subsection{Euler's Totient Theorem}

\begin{theorem}[Euler's Totient Theorem]
For $\gcd(a, n) = 1$:
\[
  a^{\phi(n)} \equiv 1 \pmod{n},
\]
where $\phi(n) = n \prod_{p \mid n}\left(1 - \tfrac{1}{p}\right)$ counts integers in $[1,n]$ coprime to $n$.
\end{theorem}

\textbf{Intuition.} Generalizes Fermat to composite moduli. The order of $a$ in $(\Z/n\Z)^*$ divides $\phi(n)$.

\textbf{Mathematical formulation.}
\[
  \phi(n) = n \prod_{p \mid n} \left(1 - \frac{1}{p}\right), \quad \phi(p^k) = p^{k-1}(p-1).
\]

\begin{cputility}
Reduces exponents modulo $\phi(n)$; useful when the modulus is not prime.
\end{cputility}

\subsubsection{Python Implementation}
\begin{lstlisting}[caption={Euler's Totient Function}]
def euler_phi(n):
    result = n
    p = 2
    temp = n
    while p * p <= temp:
        if temp % p == 0:
            while temp % p == 0:
                temp //= p
            result -= result // p
        p += 1
    if temp > 1:
        result -= result // temp
    return result

# Sieve version for range [1, N]
def phi_sieve(N):
    phi = list(range(N + 1))
    for i in range(2, N + 1):
        if phi[i] == i:          # i is prime
            for j in range(i, N + 1, i):
                phi[j] -= phi[j] // i
    return phi

print(euler_phi(12))  # 4: {1,5,7,11}
\end{lstlisting}

\textbf{Complexity.} $\bigO{\sqrt{n}}$ per value; $\bigO{n \log \log n}$ for sieve.

\begin{cpnote}
Euler's theorem requires $\gcd(a,n)=1$. If not, the formula breaks.
\end{cpnote}

% ─────────────────────────────────────────────────────────────────────────────
\subsection{Modular Inverse}

\begin{definition}
The modular inverse of $a$ modulo $m$ is $x$ such that $ax \equiv 1 \pmod{m}$. It exists iff $\gcd(a,m)=1$.
\end{definition}

\textbf{Three methods:} (1) Fermat ($m$ prime), (2) Extended GCD (general), (3) Linear sieve for range.

\subsubsection{Python Implementation}
\begin{lstlisting}[caption={Modular Inverse -- three methods}]
MOD = 10**9 + 7

# Method 1: Fermat (m must be prime)
def inv_fermat(a, m=MOD):
    return pow(a, m - 2, m)

# Method 2: Extended GCD
def ext_gcd(a, b):
    if b == 0:
        return a, 1, 0
    g, x, y = ext_gcd(b, a % b)
    return g, y, x - (a // b) * y

def inv_ext(a, m):
    g, x, _ = ext_gcd(a % m, m)
    if g != 1:
        return None  # No inverse
    return x % m

# Method 3: Linear sieve for inv[1..N]
def inv_sieve(N, m=MOD):
    inv = [0] * (N + 1)
    inv[1] = 1
    for i in range(2, N + 1):
        inv[i] = -(m // i) * inv[m % i] % m
    return inv
\end{lstlisting}

\textbf{Complexity.} Fermat: $\bigO{\log m}$; Ext GCD: $\bigO{\log m}$; Sieve: $\bigO{N}$.

% ─────────────────────────────────────────────────────────────────────────────
\subsection{Extended Euclidean Algorithm}

\begin{theorem}[Bézout's Identity]
For integers $a,b$ (not both zero), there exist integers $x,y$ such that
\[
  ax + by = \gcd(a, b).
\]
The Extended Euclidean Algorithm computes $\gcd(a,b)$, $x$, $y$ simultaneously.
\end{theorem}

\textbf{Intuition.} Backtrack the Euclidean steps to express the GCD as a linear combination.

\begin{cputility}
Solve linear Diophantine equations; compute modular inverses for non-prime moduli; CRT.
\end{cputility}

\subsubsection{Python Implementation}
\begin{lstlisting}[caption={Extended Euclidean Algorithm}]
def ext_gcd(a, b):
    """Returns (gcd, x, y) such that a*x + b*y = gcd."""
    if b == 0:
        return a, 1, 0
    g, x1, y1 = ext_gcd(b, a % b)
    return g, y1, x1 - (a // b) * y1

# Iterative version (avoids recursion limit)
def ext_gcd_iter(a, b):
    old_r, r = a, b
    old_s, s = 1, 0
    while r:
        q = old_r // r
        old_r, r = r, old_r - q * r
        old_s, s = s, old_s - q * s
    # old_s * a + ((old_r - old_s * a) // b) * b = old_r
    y = (old_r - old_s * a) // b if b else 0
    return old_r, old_s, y

g, x, y = ext_gcd(35, 15)
print(g, x, y)  # 5, 1, -2  => 35*1 + 15*(-2) = 5
\end{lstlisting}

\textbf{Complexity.} $\bigO{\log(\min(a,b))}$.

% ─────────────────────────────────────────────────────────────────────────────
\subsection{Chinese Remainder Theorem (CRT)}

\begin{theorem}[Chinese Remainder Theorem]
Let $m_1, m_2, \ldots, m_k$ be pairwise coprime moduli. Then the system
\[
  x \equiv a_i \pmod{m_i}, \quad i = 1, \ldots, k
\]
has a unique solution modulo $M = m_1 m_2 \cdots m_k$.
\end{theorem}

\textbf{Intuition.} The mapping $x \mapsto (x \bmod m_1, \ldots, x \bmod m_k)$ is a ring isomorphism $\Z/M\Z \to \Z/m_1\Z \times \cdots \times \Z/m_k\Z$.

\begin{cputility}
Combine partial results from separate moduli; reconstruct large answers; speed up DP via modular reduction.
\end{cputility}

\subsubsection{Python Implementation}
\begin{lstlisting}[caption={CRT -- pairwise coprime and general version}]
def crt(remainders, moduli):
    """CRT for pairwise coprime moduli."""
    M = 1
    for m in moduli:
        M *= m
    x = 0
    for a, m in zip(remainders, moduli):
        Mi = M // m
        _, inv, _ = ext_gcd(Mi, m)
        x += a * Mi * inv
    return x % M

def ext_gcd(a, b):
    if b == 0: return a, 1, 0
    g, x, y = ext_gcd(b, a % b)
    return g, y, x - (a // b) * y

# General CRT: merge two equations x = a1 (mod m1), x = a2 (mod m2)
def crt2(a1, m1, a2, m2):
    g, p, _ = ext_gcd(m1, m2)
    if (a2 - a1) % g != 0:
        return None, None  # No solution
    lcm = m1 // g * m2
    ans = (a1 + m1 * ((a2 - a1) // g * p % (m2 // g))) % lcm
    return ans, lcm

# Example: x=2(mod 3), x=3(mod 5), x=2(mod 7) => x=23
print(crt([2, 3, 2], [3, 5, 7]))  # 23
\end{lstlisting}

\textbf{Complexity.} $\bigO{k \log M}$.

% ─────────────────────────────────────────────────────────────────────────────
\subsection{Wilson's Theorem}

\begin{theorem}[Wilson's Theorem]
An integer $p > 1$ is prime if and only if
\[
  (p-1)! \equiv -1 \pmod{p}.
\]
\end{theorem}

\textbf{Intuition.} In $(\Z/p\Z)^*$, every element pairs with its inverse except $\pm 1$. So the product of all non-zero residues is $(-1)(1) = -1$.

\begin{cputility}
Primality test (theoretical); computing $(p-1)! \bmod p$ directly; generating problems with factorial modular structure.
\end{cputility}

\subsubsection{Python Implementation}
\begin{lstlisting}[caption={Wilson's Theorem -- primality test}]
def wilson_prime(p):
    """Returns True if p is prime (Wilson's Theorem).
    Not efficient for large p -- use as theoretical check."""
    if p < 2:
        return False
    factorial = 1
    for i in range(1, p):
        factorial = factorial * i % p
    return factorial == p - 1

# In practice: (p-1)! mod p for prime p = p-1
print(wilson_prime(7))   # True
print(wilson_prime(10))  # False
\end{lstlisting}

\textbf{Complexity.} $\bigO{p}$ -- only theoretical use, not practical for primality.

\begin{cpnote}
Wilson's theorem is rarely used directly in contests but appears in problems asking for $(n-1)! \bmod n$ or detecting primes algebraically.
\end{cpnote}

% ─────────────────────────────────────────────────────────────────────────────
\subsection{Sieve of Eratosthenes}

\begin{algorithm}[Sieve of Eratosthenes]
Mark all integers $\geq 2$ as prime. For each prime $p$ found, mark all multiples $p^2, p^2+p, \ldots$ as composite.
\end{algorithm}

\textbf{Intuition.} Any composite $n \leq N$ has a prime factor $\leq \sqrt{N}$, so crossing off multiples of primes up to $\sqrt{N}$ is sufficient.

\begin{cputility}
Precompute all primes up to $N$; smallest prime factor (SPF) for factorization; Euler's totient sieve.
\end{cputility}

\subsubsection{Python Implementation}
\begin{lstlisting}[caption={Sieve of Eratosthenes + Smallest Prime Factor sieve}]
def sieve(N):
    """Returns boolean array: is_prime[i] = True if i is prime."""
    is_prime = bytearray([1]) * (N + 1)
    is_prime[0] = is_prime[1] = 0
    for i in range(2, int(N**0.5) + 1):
        if is_prime[i]:
            is_prime[i*i::i] = bytearray(len(is_prime[i*i::i]))
    return is_prime

def spf_sieve(N):
    """Smallest prime factor for each i in [2, N]."""
    spf = list(range(N + 1))
    for i in range(2, int(N**0.5) + 1):
        if spf[i] == i:  # i is prime
            for j in range(i*i, N + 1, i):
                if spf[j] == j:
                    spf[j] = i
    return spf

def factorize_fast(n, spf):
    """Factorize n using precomputed SPF in O(log n)."""
    factors = []
    while n > 1:
        p = spf[n]
        while n % p == 0:
            factors.append(p)
            n //= p
    return factors

primes = [i for i, v in enumerate(sieve(50)) if v]
print(primes)  # [2, 3, 5, 7, 11, 13, 17, 19, 23, 29, 31, 37, 41, 43, 47]
\end{lstlisting}

\textbf{Complexity.} $\bigO{N \log \log N}$ time; $\bigO{N}$ space.

% ─────────────────────────────────────────────────────────────────────────────
\subsection{Prime Factorization}

\begin{algorithm}[Trial Division]
To factorize $n$: try all primes $p \leq \sqrt{n}$. If $p \mid n$, divide repeatedly. Any remaining factor $> 1$ is prime.
\end{algorithm}

\subsubsection{Python Implementation}
\begin{lstlisting}[caption={Prime Factorization by Trial Division}]
def factorize(n):
    """Returns dict {prime: exponent}."""
    factors = {}
    d = 2
    while d * d <= n:
        while n % d == 0:
            factors[d] = factors.get(d, 0) + 1
            n //= d
        d += 1
    if n > 1:
        factors[n] = factors.get(n, 0) + 1
    return factors

def num_divisors(n):
    f = factorize(n)
    result = 1
    for exp in f.values():
        result *= (exp + 1)
    return result

def sum_divisors(n):
    f = factorize(n)
    result = 1
    for p, e in f.items():
        result *= (p**(e+1) - 1) // (p - 1)
    return result

print(factorize(360))  # {2:3, 3:2, 5:1}
print(num_divisors(360))  # 24
\end{lstlisting}

\textbf{Complexity.} $\bigO{\sqrt{n}}$ per number; $\bigO{\log n}$ with precomputed SPF.

% ─────────────────────────────────────────────────────────────────────────────
\subsection{GCD / LCM}

\begin{theorem}[Euclidean Algorithm]
$\gcd(a, b) = \gcd(b, a \bmod b)$, with $\gcd(a, 0) = a$.
\end{theorem}

\[
  \text{lcm}(a,b) = \frac{a \cdot b}{\gcd(a,b)}.
\]

\subsubsection{Python Implementation}
\begin{lstlisting}[caption={GCD and LCM}]
from math import gcd
from functools import reduce

def lcm(a, b):
    return a // gcd(a, b) * b

def gcd_list(lst):
    return reduce(gcd, lst)

def lcm_list(lst):
    return reduce(lcm, lst)

# GCD of array
print(gcd_list([12, 18, 24]))  # 6
print(lcm_list([4, 6, 10]))    # 60
\end{lstlisting}

\textbf{Complexity.} $\bigO{\log(\min(a,b))}$.

% ─────────────────────────────────────────────────────────────────────────────
\subsection{Fast Exponentiation (Binary Exponentiation)}

\begin{algorithm}[Binary Exponentiation]
Compute $a^b \bmod m$ in $\bigO{\log b}$ by repeated squaring:
\[
  a^b = \begin{cases} 1 & b = 0 \\ (a^{b/2})^2 & b \text{ even} \\ a \cdot a^{b-1} & b \text{ odd} \end{cases}
\]
\end{algorithm}

\subsubsection{Python Implementation}
\begin{lstlisting}[caption={Binary Exponentiation -- matrix and modular versions}]
def power(base, exp, mod):
    """Compute base^exp mod (mod)."""
    result = 1
    base %= mod
    while exp > 0:
        if exp & 1:
            result = result * base % mod
        base = base * base % mod
        exp >>= 1
    return result

# Matrix exponentiation
def mat_mul(A, B, mod):
    n = len(A)
    C = [[0]*n for _ in range(n)]
    for i in range(n):
        for k in range(n):
            if A[i][k] == 0: continue
            for j in range(n):
                C[i][j] = (C[i][j] + A[i][k] * B[k][j]) % mod
    return C

def mat_pow(M, p, mod):
    n = len(M)
    result = [[int(i==j) for j in range(n)] for i in range(n)]
    while p:
        if p & 1:
            result = mat_mul(result, M, mod)
        M = mat_mul(M, M, mod)
        p >>= 1
    return result

# Fibonacci via matrix exponentiation
def fib(n, mod=10**9+7):
    if n <= 1: return n
    M = [[1,1],[1,0]]
    return mat_pow(M, n-1, mod)[0][0]

print(power(2, 100, 10**9+7))
print(fib(50))
\end{lstlisting}

\textbf{Complexity.} $\bigO{\log b}$ for scalar; $\bigO{k^3 \log b}$ for $k \times k$ matrix.

% ─────────────────────────────────────────────────────────────────────────────
\subsection{Diophantine Equations}

\begin{theorem}
The linear Diophantine equation $ax + by = c$ has integer solutions iff $\gcd(a,b) \mid c$.
If $(x_0, y_0)$ is one solution, the general solution is:
\[
  x = x_0 + \frac{b}{d}t, \quad y = y_0 - \frac{a}{d}t, \quad t \in \Z,
\]
where $d = \gcd(a,b)$.
\end{theorem}

\subsubsection{Python Implementation}
\begin{lstlisting}[caption={Linear Diophantine Equation solver}]
def ext_gcd(a, b):
    if b == 0: return a, 1, 0
    g, x, y = ext_gcd(b, a % b)
    return g, y, x - (a // b) * y

def diophantine(a, b, c):
    """Solve ax + by = c. Returns (x0, y0, step_x, step_y) or None."""
    g, x0, y0 = ext_gcd(abs(a), abs(b))
    if c % g != 0:
        return None  # No integer solution
    x0 *= c // g
    y0 *= c // g
    if a < 0: x0 = -x0
    if b < 0: y0 = -y0
    return x0, y0, b // g, a // g
    # General: x = x0 + (b/g)*t, y = y0 - (a/g)*t

sol = diophantine(3, 5, 11)
if sol:
    x0, y0, dx, dy = sol
    print(f"x={x0} + {dx}*t, y={y0} - {dy}*t")
\end{lstlisting}

\textbf{Complexity.} $\bigO{\log(\min(a,b))}$.

\begin{cpnote}
For minimizing $|x|+|y|$ or finding positive solutions, iterate $t$ to constrain the solution.
\end{cpnote}

\clearpage
% =============================================================================
\section{Combinatorics}
% =============================================================================

% ─────────────────────────────────────────────────────────────────────────────
\subsection{Binomial Theorem}

\begin{theorem}[Binomial Theorem]
For integers $n \geq 0$ and any $a, b$:
\[
  (a+b)^n = \sum_{k=0}^{n} \binom{n}{k} a^k b^{n-k}.
\]
\end{theorem}

\textbf{Intuition.} Each term in the expansion picks either $a$ or $b$ from each factor; $\binom{n}{k}$ counts the ways to choose $k$ copies of $a$.

\begin{cputility}
Expanding polynomial expressions; computing sums like $\sum_k \binom{n}{k} x^k$; generating functions.
\end{cputility}

\subsubsection{Python Implementation}
\begin{lstlisting}[caption={Binomial expansion and coefficient extraction}]
MOD = 10**9 + 7

def binomial_expansion_coeff(n, target_power_a):
    """Returns coefficient of a^k b^(n-k) in (a+b)^n."""
    k = target_power_a
    return comb_mod(n, k)

def comb_mod(n, k, mod=MOD):
    if k < 0 or k > n: return 0
    if k == 0 or k == n: return 1
    # Precompute factorials
    fact = [1] * (n + 1)
    for i in range(1, n + 1):
        fact[i] = fact[i-1] * i % mod
    inv_fact = [1] * (n + 1)
    inv_fact[n] = pow(fact[n], mod - 2, mod)
    for i in range(n - 1, -1, -1):
        inv_fact[i] = inv_fact[i+1] * (i+1) % mod
    return fact[n] * inv_fact[k] % mod * inv_fact[n-k] % mod
\end{lstlisting}

\textbf{Complexity.} $\bigO{n}$ precomputation, $\bigO{1}$ per query.

% ─────────────────────────────────────────────────────────────────────────────
\subsection{Binomial Coefficients \& Precomputed Factorials}

\[
  \binom{n}{k} = \frac{n!}{k!(n-k)!}, \quad \binom{n}{k} = \binom{n-1}{k-1} + \binom{n-1}{k}.
\]

\subsubsection{Python Implementation}
\begin{lstlisting}[caption={Factorial precomputation for fast C(n,k) mod p}]
class Combinatorics:
    def __init__(self, N, mod=10**9+7):
        self.mod = mod
        self.fact = [1] * (N + 1)
        self.inv_fact = [1] * (N + 1)
        for i in range(1, N + 1):
            self.fact[i] = self.fact[i-1] * i % mod
        self.inv_fact[N] = pow(self.fact[N], mod - 2, mod)
        for i in range(N - 1, -1, -1):
            self.inv_fact[i] = self.inv_fact[i+1] * (i+1) % mod

    def C(self, n, k):
        if k < 0 or k > n: return 0
        return self.fact[n] * self.inv_fact[k] % self.mod * self.inv_fact[n-k] % self.mod

    def P(self, n, k):
        """Permutations P(n,k) = n!/(n-k)!"""
        if k < 0 or k > n: return 0
        return self.fact[n] * self.inv_fact[n-k] % self.mod

cb = Combinatorics(200005)
print(cb.C(100, 3))   # 161700
print(cb.P(10, 3))    # 720
\end{lstlisting}

% ─────────────────────────────────────────────────────────────────────────────
\subsection{Stars and Bars}

\begin{theorem}[Stars and Bars]
The number of ways to place $n$ indistinguishable balls into $k$ distinguishable bins is:
\[
  \binom{n+k-1}{k-1}.
\]
With each bin having $\geq 1$ ball: $\binom{n-1}{k-1}$.
\end{theorem}

\textbf{Intuition.} Place $k-1$ dividers among $n+k-1$ positions.

\begin{cputility}
Counting integer solutions to $x_1+x_2+\cdots+x_k = n$ with $x_i \geq 0$.
\end{cputility}

\subsubsection{Python Implementation}
\begin{lstlisting}[caption={Stars and Bars -- non-negative and positive integer solutions}]
def stars_bars(n, k, mod=10**9+7):
    """Number of ways: x1+...+xk = n, xi >= 0."""
    return comb_mod(n + k - 1, k - 1, mod)

def stars_bars_positive(n, k, mod=10**9+7):
    """Number of ways: x1+...+xk = n, xi >= 1."""
    if n < k: return 0
    return comb_mod(n - 1, k - 1, mod)

def comb_mod(n, k, mod=10**9+7):
    if k < 0 or k > n: return 0
    # Use small pascal for small n, or factorial method
    num, den = 1, 1
    k = min(k, n - k)
    for i in range(k):
        num = num * (n - i) % mod
        den = den * (i + 1) % mod
    return num * pow(den, mod-2, mod) % mod

print(stars_bars(5, 3))   # C(7,2) = 21
\end{lstlisting}

% ─────────────────────────────────────────────────────────────────────────────
\subsection{Inclusion--Exclusion Principle}

\begin{theorem}[Inclusion--Exclusion]
For finite sets $A_1, \ldots, A_n$:
\[
  \left|\bigcup_{i=1}^n A_i\right| = \sum_{i}|A_i| - \sum_{i<j}|A_i \cap A_j| + \cdots + (-1)^{n+1}|A_1\cap\cdots\cap A_n|.
\]
\end{theorem}

\textbf{Intuition.} Add sets, subtract over-counting of pairwise intersections, add back triple intersections, etc.

\begin{cputility}
Counting elements satisfying at least one of several conditions; complementary counting; Möbius-style problems.
\end{cputility}

\subsubsection{Python Implementation}
\begin{lstlisting}[caption={Inclusion-Exclusion -- count multiples of primes in [1,N]}]
def count_multiples_iec(N, primes):
    """Count integers in [1,N] divisible by at least one prime."""
    n = len(primes)
    total = 0
    for mask in range(1, 1 << n):
        prod = 1
        bits = bin(mask).count('1')
        for i in range(n):
            if mask & (1 << i):
                prod *= primes[i]
                if prod > N: break
        else:
            cnt = N // prod
            if bits % 2 == 1:
                total += cnt
            else:
                total -= cnt
    return total

print(count_multiples_iec(30, [2, 3, 5]))  # 22
\end{lstlisting}

\textbf{Complexity.} $\bigO{2^n}$ for $n$ sets.

% ─────────────────────────────────────────────────────────────────────────────
\subsection{Pigeonhole Principle}

\begin{theorem}[Pigeonhole Principle]
If $n > k$ objects are distributed into $k$ containers, at least one container has $\geq 2$ objects. More generally, some container has $\geq \lceil n/k \rceil$ objects.
\end{theorem}

\textbf{Intuition.} Averaging argument: if all bins had $< \lceil n/k \rceil$, the total would be less than $n$.

\begin{cputility}
Proving existence; bounding maximum; guarantee-type arguments; hash collisions.
\end{cputility}

\textbf{Example.} In any sequence of $n^2+1$ distinct integers, there is a monotone subsequence of length $n+1$ (Dilworth / Erdős–Szekeres).

% ─────────────────────────────────────────────────────────────────────────────
\subsection{Catalan Numbers}

\begin{theorem}[Catalan Numbers]
The $n$-th Catalan number is:
\[
  C_n = \frac{1}{n+1}\binom{2n}{n} = \frac{(2n)!}{(n+1)!\,n!}.
\]
$C_0=1, C_1=1, C_2=2, C_3=5, C_4=14, \ldots$
Recurrence: $C_{n+1} = \sum_{i=0}^n C_i C_{n-i}$.
\end{theorem}

\textbf{Where to use.} Balanced parentheses, BST shapes, triangulations, non-crossing partitions, mountain ranges (Dyck paths).

\subsubsection{Python Implementation}
\begin{lstlisting}[caption={Catalan Numbers mod p}]
def catalan(n, mod=10**9+7):
    """C_n = C(2n,n)/(n+1) mod p."""
    # Precompute factorials
    N = 2 * n + 2
    fact = [1] * (N + 1)
    for i in range(1, N + 1):
        fact[i] = fact[i-1] * i % mod
    inv_fact = [1] * (N + 1)
    inv_fact[N] = pow(fact[N], mod-2, mod)
    for i in range(N-1, -1, -1):
        inv_fact[i] = inv_fact[i+1] * (i+1) % mod
    return fact[2*n] * inv_fact[n] % mod * inv_fact[n+1] % mod

# First 10 Catalan numbers
for i in range(10):
    print(i, catalan(i))
# 0:1, 1:1, 2:2, 3:5, 4:14, 5:42, 6:132, 7:429, 8:1430, 9:4862
\end{lstlisting}

\textbf{Complexity.} $\bigO{n}$ per value.

% ─────────────────────────────────────────────────────────────────────────────
\subsection{Sperner's Theorem}

\begin{theorem}[Sperner's Theorem]
The maximum size of an antichain (a family of sets no one of which contains another) in the power set of an $n$-element set is $\binom{n}{\lfloor n/2 \rfloor}$.
\end{theorem}

\textbf{Intuition.} Among all layers of the Boolean lattice $\binom{n}{0}, \binom{n}{1}, \ldots, \binom{n}{n}$, the middle layer $\binom{n}{\lfloor n/2 \rfloor}$ is the largest, and no two sets of equal size can contain one another.

\textbf{Mathematical formulation (LYM inequality).}
\[
  \sum_{A \in \mathcal{F}} \frac{1}{\binom{n}{|A|}} \leq 1 \quad \Longrightarrow \quad |\mathcal{F}| \leq \binom{n}{\lfloor n/2 \rfloor}.
\]

\begin{cputility}
Bounding families of subsets; proving impossibility of certain constructions; extremal set-system arguments.
\end{cputility}

\textbf{Where to use.} Problems asking for maximum subset families with no containment; proof arguments in combinatorial optimization.

\textbf{Example Problem.} Given $n$ elements, find the maximum number of subsets such that no subset is contained in another.

\subsubsection{Python Implementation}
\begin{lstlisting}[caption={Sperner's Theorem -- max antichain and verification}]
from math import comb

def max_antichain_size(n):
    """Maximum antichain in power set of {1,...,n}."""
    return comb(n, n // 2)

def is_antichain(family):
    """Check if a family of sets forms an antichain."""
    fam = [set(s) for s in family]
    for i in range(len(fam)):
        for j in range(len(fam)):
            if i != j and fam[i] <= fam[j]:
                return False
    return True

# For n=4, max antichain = C(4,2) = 6 sets
print(max_antichain_size(4))  # 6
family = [{0,1},{0,2},{0,3},{1,2},{1,3},{2,3}]
print(is_antichain(family))   # True
\end{lstlisting}

\textbf{Complexity.} Computing $\binom{n}{\lfloor n/2\rfloor}$: $\bigO{n}$; verifying antichain: $\bigO{|\mathcal{F}|^2 \cdot n}$.

% ─────────────────────────────────────────────────────────────────────────────
\subsection{Derangements}

\begin{definition}
A derangement of $n$ elements is a permutation with no fixed points. The count is:
\[
  D_n = n!\sum_{k=0}^{n}\frac{(-1)^k}{k!} \approx \frac{n!}{e}.
\]
Recurrence: $D_n = (n-1)(D_{n-1} + D_{n-2})$, $D_0=1, D_1=0$.
\end{definition}

\subsubsection{Python Implementation}
\begin{lstlisting}[caption={Derangements mod p}]
def derangements(N, mod=10**9+7):
    """Returns D[0..N] where D[n] = number of derangements of n."""
    D = [0] * (N + 1)
    D[0] = 1
    if N >= 1: D[1] = 0
    for n in range(2, N + 1):
        D[n] = (n - 1) * (D[n-1] + D[n-2]) % mod
    return D

D = derangements(10)
print(D[4])  # 9
print(D[5])  # 44
\end{lstlisting}

% ─────────────────────────────────────────────────────────────────────────────
\subsection{Combinations with Repetition}

\begin{theorem}
The number of multisets of size $k$ from an $n$-element set is:
\[
  \binom{n+k-1}{k} = \binom{n+k-1}{n-1}.
\]
\end{theorem}

\textbf{Intuition.} Equivalent to Stars and Bars: choose $k$ elements with repetition from $n$ types.

\subsubsection{Python Implementation}
\begin{lstlisting}[caption={Combinations with repetition}]
def comb_rep(n, k, mod=10**9+7):
    """Multisets: choose k from n types with repetition."""
    return comb_mod(n + k - 1, k, mod)

def comb_mod(n, k, mod=10**9+7):
    if k < 0 or k > n: return 0
    k = min(k, n-k)
    num = den = 1
    for i in range(k):
        num = num * (n - i) % mod
        den = den * (i + 1) % mod
    return num * pow(den, mod-2, mod) % mod

print(comb_rep(3, 2))  # 6 (aa,ab,ac,bb,bc,cc)
\end{lstlisting}

\clearpage
% =============================================================================
\section{Graph Theory}
% =============================================================================

\subsection{BFS -- Breadth-First Search}

\begin{algorithm}[BFS]
Explore vertices level by level from a source. Uses a queue. Finds shortest paths in unweighted graphs.
\end{algorithm}

\textbf{Where to use.} Shortest path in unweighted graph; 0-1 BFS; level-order traversal; bipartite check.

\textbf{Typical CP scenarios.} Grid shortest path (4/8-directional); multi-source BFS (spread from multiple sources simultaneously); BFS on states (e.g., Rubik's cube, game states).

\textbf{Example Problem.} Given a grid with \texttt{\#} walls and \texttt{.} paths, find the shortest path from top-left to bottom-right.

\subsubsection{Python Implementation}
\begin{lstlisting}[caption={BFS -- shortest paths}]
from collections import deque

def bfs(graph, src, n):
    """graph: adjacency list. Returns dist[] from src."""
    dist = [-1] * n
    dist[src] = 0
    q = deque([src])
    while q:
        u = q.popleft()
        for v in graph[u]:
            if dist[v] == -1:
                dist[v] = dist[u] + 1
                q.append(v)
    return dist

# 0-1 BFS for graphs with edge weights 0 or 1
def bfs_01(graph, src, n):
    dist = [float('inf')] * n
    dist[src] = 0
    dq = deque([src])
    while dq:
        u = dq.popleft()
        for v, w in graph[u]:
            if dist[u] + w < dist[v]:
                dist[v] = dist[u] + w
                if w == 0:
                    dq.appendleft(v)
                else:
                    dq.append(v)
    return dist
\end{lstlisting}

\textbf{Complexity.} $\bigO{V + E}$.

\begin{cpnote}
For multi-source BFS, push all sources into the queue initially with distance 0. For states involving $\{\text{position}, \text{keys collected}\}$, BFS on the augmented state space.
\end{cpnote}

% ─────────────────────────────────────────────────────────────────────────────
\subsection{DFS -- Depth-First Search}

\begin{algorithm}[DFS]
Explore as deep as possible before backtracking. Can be recursive or iterative. Produces DFS tree, discovery/finish times.
\end{algorithm}

\textbf{Where to use.} Connected components; cycle detection; topological sort; SCC; bridges/articulation points; backtracking.

\textbf{Typical CP scenarios.} Finding connected components; detecting back edges (cycles); tree diameter via two DFS calls; flood fill.

\textbf{Example Problem.} Given an undirected graph, count the number of connected components.

\subsubsection{Python Implementation}
\begin{lstlisting}[caption={DFS -- iterative with timestamps}]
def dfs(graph, n):
    visited = [False] * n
    order = []
    def dfs_visit(u):
        stack = [(u, 0)]
        while stack:
            v, idx = stack.pop()
            if idx == 0:
                if visited[v]: continue
                visited[v] = True
                order.append(v)
                stack.append((v, 1))
                for w in reversed(graph[v]):
                    if not visited[w]:
                        stack.append((w, 0))
    for i in range(n):
        if not visited[i]:
            dfs_visit(i)
    return order
\end{lstlisting}

\textbf{Complexity.} $\bigO{V + E}$.

\begin{cpnote}
Python's default recursion limit is 1000. For DFS on large graphs, always use iterative DFS or raise the limit with \texttt{sys.setrecursionlimit()} and use threading for larger stack.
\end{cpnote}

% ─────────────────────────────────────────────────────────────────────────────
\subsection{Dijkstra's Algorithm}

\begin{algorithm}[Dijkstra]
Single-source shortest paths in graphs with non-negative edge weights. Uses a priority queue (min-heap).
\end{algorithm}

\textbf{Where to use.} Weighted graphs with non-negative edges; network routing; shortest path in state spaces (BFS with costs).

\textbf{Typical CP scenarios.} Road networks; weighted grids; shortest path with constraints (multi-state Dijkstra).

\textbf{Example Problem.} Given $n$ cities and $m$ weighted roads, find the shortest distance from city 1 to city $n$.

\[
  d[v] = \min_{(u,v)\in E} (d[u] + w(u,v))
\]

\subsubsection{Python Implementation}
\begin{lstlisting}[caption={Dijkstra with heapq}]
import heapq

def dijkstra(graph, src, n):
    """graph[u] = [(v, w), ...]. Returns dist[] from src."""
    dist = [float('inf')] * n
    dist[src] = 0
    heap = [(0, src)]
    while heap:
        d, u = heapq.heappop(heap)
        if d > dist[u]:
            continue
        for v, w in graph[u]:
            if dist[u] + w < dist[v]:
                dist[v] = dist[u] + w
                heapq.heappush(heap, (dist[v], v))
    return dist

# Example
n = 5
g = [[] for _ in range(n)]
g[0].append((1, 2)); g[0].append((2, 4))
g[1].append((2, 1)); g[1].append((3, 7))
g[2].append((4, 3)); g[3].append((4, 1))
print(dijkstra(g, 0, n))  # [0, 2, 3, 9, 6]
\end{lstlisting}

\textbf{Complexity.} $\bigO{(V+E)\log V}$ with binary heap.

\begin{cpnote}
Dijkstra fails with negative edges. Use Bellman-Ford instead.
\end{cpnote}

% ─────────────────────────────────────────────────────────────────────────────
\subsection{Bellman-Ford Algorithm}

\begin{algorithm}[Bellman-Ford]
Relax all edges $V-1$ times. Detects negative cycles in a $V$-th relaxation pass.
\end{algorithm}

\textbf{Where to use.} Graphs with negative edge weights; detecting negative cycles; when Dijkstra is inapplicable.

\textbf{Typical CP scenarios.} Currency arbitrage (negative cycle = profit); shortest path with negative edges; SPFA optimization.

\textbf{Example Problem.} Given a currency exchange graph, determine if there is an arbitrage opportunity (negative cycle).

\subsubsection{Python Implementation}
\begin{lstlisting}[caption={Bellman-Ford with negative cycle detection}]
def bellman_ford(edges, src, n):
    """edges = [(u,v,w)]. Returns (dist, has_neg_cycle)."""
    dist = [float('inf')] * n
    dist[src] = 0
    for _ in range(n - 1):
        for u, v, w in edges:
            if dist[u] + w < dist[v]:
                dist[v] = dist[u] + w
    # Check for negative cycle
    neg_cycle = False
    for u, v, w in edges:
        if dist[u] + w < dist[v]:
            neg_cycle = True
            break
    return dist, neg_cycle
\end{lstlisting}

\textbf{Complexity.} $\bigO{VE}$.

\begin{cpnote}
SPFA (Shortest Path Faster Algorithm) is Bellman-Ford with a queue -- often much faster in practice but $\bigO{VE}$ worst case. Avoid on adversarial test cases.
\end{cpnote}

% ─────────────────────────────────────────────────────────────────────────────
\subsection{Floyd-Warshall Algorithm}

\begin{algorithm}[Floyd-Warshall]
All-pairs shortest paths in $\bigO{V^3}$.
\[
  d[i][j] = \min(d[i][j],\; d[i][k] + d[k][j]) \quad \forall k
\]
\end{algorithm}

\textbf{Where to use.} Small dense graphs ($V \leq 500$); transitive closure; detecting negative cycles ($d[i][i] < 0$).

\textbf{Typical CP scenarios.} All-pairs shortest paths; reachability; minimum cycle through a node.

\textbf{Example Problem.} Given $n \leq 400$ cities with direct flights, find the shortest round trip through any city.

\subsubsection{Python Implementation}
\begin{lstlisting}[caption={Floyd-Warshall all-pairs shortest paths}]
INF = float('inf')

def floyd_warshall(dist):
    """dist[i][j] = initial edge weight (INF if no edge), in-place."""
    n = len(dist)
    for k in range(n):
        for i in range(n):
            if dist[i][k] == INF: continue
            for j in range(n):
                if dist[i][k] + dist[k][j] < dist[i][j]:
                    dist[i][j] = dist[i][k] + dist[k][j]
    # Negative cycle: dist[i][i] < 0
    return dist

# Usage
n = 4
d = [[INF]*n for _ in range(n)]
for i in range(n): d[i][i] = 0
d[0][1]=3; d[0][3]=7; d[1][0]=8; d[1][2]=2
d[2][0]=5; d[2][3]=1; d[3][0]=2
floyd_warshall(d)
\end{lstlisting}

\textbf{Complexity.} $\bigO{V^3}$. Best for dense graphs with $V \leq 500$.

% ─────────────────────────────────────────────────────────────────────────────
\subsection{Topological Sort}

\begin{algorithm}[Topological Sort (Kahn's + DFS)]
Linear ordering of vertices in a DAG such that for every edge $(u,v)$, $u$ comes before $v$.
\end{algorithm}

\subsubsection{Python Implementation}
\begin{lstlisting}[caption={Topological Sort -- Kahn's BFS and DFS}]
from collections import deque

def topo_kahn(graph, n):
    """Kahn's algorithm. Returns [] if cycle exists."""
    indegree = [0] * n
    for u in range(n):
        for v in graph[u]:
            indegree[v] += 1
    q = deque(i for i in range(n) if indegree[i] == 0)
    order = []
    while q:
        u = q.popleft()
        order.append(u)
        for v in graph[u]:
            indegree[v] -= 1
            if indegree[v] == 0:
                q.append(v)
    return order if len(order) == n else []  # Empty = cycle

def topo_dfs(graph, n):
    """DFS-based topological sort."""
    visited = [0] * n  # 0=unvisited, 1=in stack, 2=done
    order = []
    def dfs(u):
        visited[u] = 1
        for v in graph[u]:
            if visited[v] == 1: return False  # Cycle
            if visited[v] == 0:
                if not dfs(v): return False
        visited[u] = 2
        order.append(u)
        return True
    for i in range(n):
        if visited[i] == 0:
            if not dfs(i): return []
    return order[::-1]
\end{lstlisting}

\textbf{Complexity.} $\bigO{V+E}$.

% ─────────────────────────────────────────────────────────────────────────────
\subsection{Kruskal's Algorithm (MST)}

\begin{algorithm}[Kruskal]
Sort edges by weight; greedily add the minimum weight edge that doesn't form a cycle (use DSU).
\end{algorithm}

\subsubsection{Python Implementation}
\begin{lstlisting}[caption={Kruskal's MST with DSU}]
class DSU:
    def __init__(self, n):
        self.p = list(range(n))
        self.rank = [0] * n

    def find(self, x):
        while self.p[x] != x:
            self.p[x] = self.p[self.p[x]]
            x = self.p[x]
        return x

    def union(self, a, b):
        a, b = self.find(a), self.find(b)
        if a == b: return False
        if self.rank[a] < self.rank[b]: a, b = b, a
        self.p[b] = a
        if self.rank[a] == self.rank[b]: self.rank[a] += 1
        return True

def kruskal(n, edges):
    """edges = [(w,u,v)]. Returns (mst_weight, mst_edges)."""
    edges.sort()
    dsu = DSU(n)
    mst_w = 0
    mst_edges = []
    for w, u, v in edges:
        if dsu.union(u, v):
            mst_w += w
            mst_edges.append((u, v, w))
            if len(mst_edges) == n - 1:
                break
    return mst_w, mst_edges
\end{lstlisting}

\textbf{Complexity.} $\bigO{E \log E}$.

% ─────────────────────────────────────────────────────────────────────────────
\subsection{Prim's Algorithm (MST)}

\subsubsection{Python Implementation}
\begin{lstlisting}[caption={Prim's MST with priority queue}]
import heapq

def prim(graph, n, src=0):
    """graph[u] = [(v,w)]. Returns MST weight."""
    visited = [False] * n
    heap = [(0, src)]
    total = 0
    while heap:
        w, u = heapq.heappop(heap)
        if visited[u]: continue
        visited[u] = True
        total += w
        for v, wt in graph[u]:
            if not visited[v]:
                heapq.heappush(heap, (wt, v))
    return total
\end{lstlisting}

\textbf{Complexity.} $\bigO{(V+E)\log V}$.

% ─────────────────────────────────────────────────────────────────────────────
\subsection{Union-Find (DSU)}

Already shown in Kruskal; additional features below.

\subsubsection{Python Implementation}
\begin{lstlisting}[caption={DSU with size tracking and rollback}]
class DSU:
    def __init__(self, n):
        self.p = list(range(n))
        self.size = [1] * n
        self.n_components = n

    def find(self, x):
        r = x
        while self.p[r] != r: r = self.p[r]
        while self.p[x] != r:
            self.p[x], x = r, self.p[x]
        return r

    def union(self, a, b):
        a, b = self.find(a), self.find(b)
        if a == b: return False
        if self.size[a] < self.size[b]: a, b = b, a
        self.p[b] = a
        self.size[a] += self.size[b]
        self.n_components -= 1
        return True

    def same(self, a, b): return self.find(a) == self.find(b)
    def comp_size(self, a): return self.size[self.find(a)]
\end{lstlisting}

\textbf{Complexity.} Near $\bigO{1}$ amortized (inverse Ackermann).

% ─────────────────────────────────────────────────────────────────────────────
\subsection{Tarjan's SCC}

\begin{algorithm}[Tarjan SCC]
DFS-based; maintains a stack and discovery/low-link times. SCCs are popped when root is identified.
\end{algorithm}

\subsubsection{Python Implementation}
\begin{lstlisting}[caption={Tarjan's Strongly Connected Components}]
import sys
sys.setrecursionlimit(300000)

def tarjan_scc(graph, n):
    disc = [-1] * n
    low = [0] * n
    on_stack = [False] * n
    stack = []
    sccs = []
    timer = [0]

    def dfs(u):
        disc[u] = low[u] = timer[0]
        timer[0] += 1
        stack.append(u)
        on_stack[u] = True
        for v in graph[u]:
            if disc[v] == -1:
                dfs(v)
                low[u] = min(low[u], low[v])
            elif on_stack[v]:
                low[u] = min(low[u], disc[v])
        if low[u] == disc[u]:
            scc = []
            while True:
                v = stack.pop()
                on_stack[v] = False
                scc.append(v)
                if v == u: break
            sccs.append(scc)

    for i in range(n):
        if disc[i] == -1:
            dfs(i)
    return sccs
\end{lstlisting}

\textbf{Complexity.} $\bigO{V+E}$.

\begin{cpnote}
This recursive implementation hits Python's recursion limit for $n > 10^4$. For large graphs, use Kosaraju's (iterative) or implement an iterative Tarjan using an explicit stack.
\end{cpnote}

% ─────────────────────────────────────────────────────────────────────────────
\subsection{Kosaraju's SCC}

\subsubsection{Python Implementation}
\begin{lstlisting}[caption={Kosaraju's SCC}]
def kosaraju(graph, rgraph, n):
    """graph and rgraph: adjacency lists."""
    visited = [False] * n
    order = []

    def dfs1(u):
        stack = [(u, 0)]
        while stack:
            v, state = stack.pop()
            if state == 0:
                if visited[v]: continue
                visited[v] = True
                stack.append((v, 1))
                for w in graph[v]:
                    if not visited[w]:
                        stack.append((w, 0))
            else:
                order.append(v)

    for i in range(n):
        if not visited[i]:
            dfs1(i)

    comp = [-1] * n
    def dfs2(u, c):
        stack = [u]
        while stack:
            v = stack.pop()
            if comp[v] != -1: continue
            comp[v] = c
            for w in rgraph[v]:
                if comp[w] == -1:
                    stack.append(w)

    c = 0
    for u in reversed(order):
        if comp[u] == -1:
            dfs2(u, c)
            c += 1
    return comp, c
\end{lstlisting}

\textbf{Complexity.} $\bigO{V+E}$.

% ─────────────────────────────────────────────────────────────────────────────
\subsection{Bridges and Articulation Points}

\subsubsection{Python Implementation}
\begin{lstlisting}[caption={Bridges and Articulation Points via DFS}]
def find_bridges_ap(graph, n):
    disc = [-1] * n
    low = [0] * n
    parent = [-1] * n
    bridges = []
    ap = [False] * n
    timer = [0]

    def dfs(u):
        disc[u] = low[u] = timer[0]
        timer[0] += 1
        child_count = 0
        for v in graph[u]:
            if disc[v] == -1:
                child_count += 1
                parent[v] = u
                dfs(v)
                low[u] = min(low[u], low[v])
                if low[v] > disc[u]:
                    bridges.append((u, v))
                if parent[u] == -1 and child_count > 1:
                    ap[u] = True
                if parent[u] != -1 and low[v] >= disc[u]:
                    ap[u] = True
            elif v != parent[u]:
                low[u] = min(low[u], disc[v])

    for i in range(n):
        if disc[i] == -1:
            dfs(i)
    return bridges, [i for i in range(n) if ap[i]]
\end{lstlisting}

\textbf{Complexity.} $\bigO{V+E}$.

\begin{cpnote}
This recursive implementation may exceed Python's stack limit for $n > 10^4$. Use \texttt{sys.setrecursionlimit} or rewrite iteratively. For multi-edge graphs, track edge index rather than \texttt{v != parent[u]} to handle duplicate edges correctly.
\end{cpnote}

% ─────────────────────────────────────────────────────────────────────────────
\subsection{LCA with Binary Lifting}

\subsubsection{Python Implementation}
\begin{lstlisting}[caption={LCA using Binary Lifting}]
import math

class LCA:
    def __init__(self, tree, root, n):
        LOG = max(1, math.floor(math.log2(n)) + 1) if n > 1 else 1
        self.LOG = LOG
        self.depth = [0] * n
        self.up = [[-1] * n for _ in range(LOG)]
        # BFS to set depth and up[0]
        from collections import deque
        visited = [False] * n
        q = deque([root])
        visited[root] = True
        self.up[0][root] = root
        while q:
            u = q.popleft()
            for v in tree[u]:
                if not visited[v]:
                    visited[v] = True
                    self.depth[v] = self.depth[u] + 1
                    self.up[0][v] = u
                    q.append(v)
        for k in range(1, LOG):
            for v in range(n):
                self.up[k][v] = self.up[k-1][self.up[k-1][v]]

    def lca(self, u, v):
        if self.depth[u] < self.depth[v]:
            u, v = v, u
        diff = self.depth[u] - self.depth[v]
        for k in range(self.LOG):
            if (diff >> k) & 1:
                u = self.up[k][u]
        if u == v: return u
        for k in range(self.LOG - 1, -1, -1):
            if self.up[k][u] != self.up[k][v]:
                u = self.up[k][u]
                v = self.up[k][v]
        return self.up[0][u]
\end{lstlisting}

\textbf{Complexity.} $\bigO{N \log N}$ build; $\bigO{\log N}$ query.

% ─────────────────────────────────────────────────────────────────────────────
\subsection{Euler Tour}

\subsubsection{Python Implementation}
\begin{lstlisting}[caption={Euler Tour -- in/out times for subtree queries}]
def euler_tour(tree, root, n):
    tin = [0] * n
    tout = [0] * n
    order = []
    timer = [0]
    stack = [(root, -1, False)]
    while stack:
        u, par, leaving = stack.pop()
        if leaving:
            tout[u] = timer[0]
            timer[0] += 1
        else:
            tin[u] = timer[0]
            timer[0] += 1
            order.append(u)
            stack.append((u, par, True))
            for v in tree[u]:
                if v != par:
                    stack.append((v, u, False))
    return tin, tout, order

# Node v is in subtree of u iff tin[u] <= tin[v] < tout[u]
\end{lstlisting}

\textbf{Complexity.} $\bigO{N}$.

% ─────────────────────────────────────────────────────────────────────────────
\subsection{Maximum Flow -- Dinic's Algorithm}

\subsubsection{Python Implementation}
\begin{lstlisting}[caption={Dinic's Maximum Flow}]
from collections import deque

class Dinic:
    def __init__(self, n):
        self.n = n
        self.graph = [[] for _ in range(n)]

    def add_edge(self, u, v, cap):
        self.graph[u].append([v, cap, len(self.graph[v])])
        self.graph[v].append([u, 0, len(self.graph[u])-1])

    def bfs(self, s, t):
        self.level = [-1] * self.n
        self.level[s] = 0
        q = deque([s])
        while q:
            u = q.popleft()
            for v, cap, _ in self.graph[u]:
                if cap > 0 and self.level[v] < 0:
                    self.level[v] = self.level[u] + 1
                    q.append(v)
        return self.level[t] >= 0

    def dfs(self, u, t, f):
        if u == t: return f
        while self.iter[u] < len(self.graph[u]):
            v, cap, rev = self.graph[u][self.iter[u]]
            if cap > 0 and self.level[v] == self.level[u] + 1:
                d = self.dfs(v, t, min(f, cap))
                if d > 0:
                    self.graph[u][self.iter[u]][1] -= d
                    self.graph[v][rev][1] += d
                    return d
            self.iter[u] += 1
        return 0

    def max_flow(self, s, t):
        flow = 0
        while self.bfs(s, t):
            self.iter = [0] * self.n
            while True:
                f = self.dfs(s, t, float('inf'))
                if f == 0: break
                flow += f
        return flow
\end{lstlisting}

\textbf{Complexity.} $\bigO{V^2 E}$; $\bigO{E\sqrt{V}}$ for unit graphs.

% ─────────────────────────────────────────────────────────────────────────────
\subsection{Bipartite Matching (Hopcroft-Karp)}

\subsubsection{Python Implementation}
\begin{lstlisting}[caption={Hopcroft-Karp maximum bipartite matching}]
from collections import deque
INF = float('inf')

def hopcroft_karp(adj, n, m):
    """adj[u] = list of v for left u, right v. Returns matching size."""
    match_l = [-1] * n
    match_r = [-1] * m
    dist = [0] * n

    def bfs():
        queue = deque()
        for u in range(n):
            if match_l[u] == -1:
                dist[u] = 0
                queue.append(u)
            else:
                dist[u] = INF
        found = False
        while queue:
            u = queue.popleft()
            for v in adj[u]:
                w = match_r[v]
                if w == -1:
                    found = True
                elif dist[w] == INF:
                    dist[w] = dist[u] + 1
                    queue.append(w)
        return found

    def dfs(u):
        for v in adj[u]:
            w = match_r[v]
            if w == -1 or (dist[w] == dist[u] + 1 and dfs(w)):
                match_l[u] = v
                match_r[v] = u
                return True
        dist[u] = INF
        return False

    matching = 0
    while bfs():
        for u in range(n):
            if match_l[u] == -1 and dfs(u):
                matching += 1
    return matching, match_l, match_r
\end{lstlisting}

\textbf{Complexity.} $\bigO{E\sqrt{V}}$.

% ─────────────────────────────────────────────────────────────────────────────
\subsection{Minimum Cut}

\begin{theorem}[Max-Flow Min-Cut Theorem]
In a flow network, the maximum flow equals the minimum cut capacity.
\end{theorem}

\textbf{Finding the min-cut.} After running max flow, do BFS/DFS from source in the residual graph. Edges from reachable to non-reachable vertices form the min cut.

\textbf{Example Problem.} Given a network of pipes with capacities, find the minimum total capacity of pipes to cut so no water flows from source to sink.

\subsubsection{Python Implementation}
\begin{lstlisting}[caption={Minimum Cut -- extract cut edges after max flow}]
from collections import deque

def find_min_cut(dinic, s):
    """After dinic.max_flow(s, t), find reachable set from s
    in residual graph. Cut = edges from reachable to rest."""
    reachable = set()
    q = deque([s])
    reachable.add(s)
    while q:
        u = q.popleft()
        for v, cap, _ in dinic.graph[u]:
            if cap > 0 and v not in reachable:
                reachable.add(v)
                q.append(v)
    # Cut edges: u reachable, v not, original cap > 0
    cut = []
    for u in reachable:
        for v, cap, rev in dinic.graph[u]:
            if v not in reachable and cap == 0:
                orig = dinic.graph[v][rev][1]
                if orig > 0:
                    cut.append((u, v, orig))
    return reachable, cut

# Usage: flow = dinic.max_flow(s, t)
# reachable, cut_edges = find_min_cut(dinic, s)
# sum of cut capacities == flow
\end{lstlisting}

\textbf{Complexity.} BFS is $\bigO{V+E}$ on top of the max-flow computation.

\begin{cpnote}
For minimum vertex cut, split each vertex $v$ into $v_{\text{in}}$ and $v_{\text{out}}$ with an edge of capacity 1, then run max flow.
\end{cpnote}

% ─────────────────────────────────────────────────────────────────────────────
\subsection{Bipartite Graph Theorem}

\begin{theorem}
A graph is bipartite if and only if it contains no odd-length cycle.
\end{theorem}

\textbf{Test via 2-coloring BFS/DFS.}

\subsubsection{Python Implementation}
\begin{lstlisting}[caption={Bipartite check via 2-coloring}]
from collections import deque

def is_bipartite(graph, n):
    color = [-1] * n
    for start in range(n):
        if color[start] != -1: continue
        color[start] = 0
        q = deque([start])
        while q:
            u = q.popleft()
            for v in graph[u]:
                if color[v] == -1:
                    color[v] = 1 - color[u]
                    q.append(v)
                elif color[v] == color[u]:
                    return False, color
    return True, color
\end{lstlisting}

\clearpage
% =============================================================================
\section{Dynamic Programming}
% =============================================================================

\subsection{0/1 Knapsack}

\begin{algorithm}[0/1 Knapsack]
Given items with weights $w_i$ and values $v_i$, maximize $\sum v_i x_i$ subject to $\sum w_i x_i \leq W$, $x_i \in \{0,1\}$.
\[
  dp[j] = \max(dp[j],\; dp[j - w_i] + v_i) \quad \text{(iterate } j \text{ backwards)}
\]
\end{algorithm}

\subsubsection{Python Implementation}
\begin{lstlisting}[caption={0/1 Knapsack}]
def knapsack_01(items, W):
    """items = [(weight, value)]. Returns max value."""
    dp = [0] * (W + 1)
    for w, v in items:
        for j in range(W, w - 1, -1):
            dp[j] = max(dp[j], dp[j - w] + v)
    return dp[W]

items = [(2,3),(3,4),(4,5),(5,6)]
print(knapsack_01(items, 8))  # 10
\end{lstlisting}

\textbf{Complexity.} $\bigO{nW}$.

% ─────────────────────────────────────────────────────────────────────────────
\subsection{Unbounded Knapsack}

\subsubsection{Python Implementation}
\begin{lstlisting}[caption={Unbounded Knapsack}]
def knapsack_unbounded(items, W):
    dp = [0] * (W + 1)
    for j in range(1, W + 1):
        for w, v in items:
            if w <= j:
                dp[j] = max(dp[j], dp[j - w] + v)
    return dp[W]
\end{lstlisting}

\textbf{Key difference from 0/1.} Iterate $j$ \emph{forward} (each item usable multiple times).

% ─────────────────────────────────────────────────────────────────────────────
\subsection{Coin Change}

\subsubsection{Python Implementation}
\begin{lstlisting}[caption={Coin Change -- min coins and count ways}]
def coin_change_min(coins, amount):
    """Min number of coins to make amount. -1 if impossible."""
    dp = [float('inf')] * (amount + 1)
    dp[0] = 0
    for c in coins:
        for j in range(c, amount + 1):
            dp[j] = min(dp[j], dp[j - c] + 1)
    return dp[amount] if dp[amount] != float('inf') else -1

def coin_change_ways(coins, amount, mod=10**9+7):
    """Count distinct ways to make amount."""
    dp = [0] * (amount + 1)
    dp[0] = 1
    for c in coins:
        for j in range(c, amount + 1):
            dp[j] = (dp[j] + dp[j - c]) % mod
    return dp[amount]

print(coin_change_min([1,5,6,9], 11))  # 2 (5+6)
print(coin_change_ways([1,2,5], 5))    # 4
\end{lstlisting}

% ─────────────────────────────────────────────────────────────────────────────
\subsection{LIS -- Longest Increasing Subsequence $\bigO{n \log n}$}

\begin{algorithm}[LIS in $\bigO{n \log n}$]
Maintain array $\mathtt{tails}$ where $\mathtt{tails}[i]$ is the smallest tail of all increasing subsequences of length $i+1$. Use binary search to update.
\end{algorithm}

\subsubsection{Python Implementation}
\begin{lstlisting}[caption={LIS O(n log n) with patience sorting}]
import bisect

def lis(arr):
    """Returns length of LIS."""
    tails = []
    for x in arr:
        pos = bisect.bisect_left(tails, x)
        if pos == len(tails):
            tails.append(x)
        else:
            tails[pos] = x
    return len(tails)

def lis_sequence(arr):
    """Returns one actual LIS (not just the length)."""
    n = len(arr)
    tails = []       # tails[i] = smallest tail of IS length i+1
    tail_idx = []     # index in arr for tails[i]
    prev_idx = [-1] * n
    for i, x in enumerate(arr):
        p = bisect.bisect_left(tails, x)
        if p == len(tails):
            tails.append(x)
            tail_idx.append(i)
        else:
            tails[p] = x
            tail_idx[p] = i
        prev_idx[i] = tail_idx[p - 1] if p > 0 else -1
    # Reconstruct backwards from last element
    result = []
    idx = tail_idx[-1]
    while idx != -1:
        result.append(arr[idx])
        idx = prev_idx[idx]
    return result[::-1]

print(lis([10,9,2,5,3,7,101,18]))  # 4
print(lis_sequence([10,9,2,5,3,7,101,18]))  # [2,3,7,18]
\end{lstlisting}

\textbf{Complexity.} $\bigO{n \log n}$.

% ─────────────────────────────────────────────────────────────────────────────
\subsection{LCS -- Longest Common Subsequence}

\subsubsection{Python Implementation}
\begin{lstlisting}[caption={LCS -- O(nm) DP}]
def lcs(s, t):
    n, m = len(s), len(t)
    dp = [[0]*(m+1) for _ in range(n+1)]
    for i in range(1, n+1):
        for j in range(1, m+1):
            if s[i-1] == t[j-1]:
                dp[i][j] = dp[i-1][j-1] + 1
            else:
                dp[i][j] = max(dp[i-1][j], dp[i][j-1])
    # Reconstruct
    seq = []
    i, j = n, m
    while i > 0 and j > 0:
        if s[i-1] == t[j-1]:
            seq.append(s[i-1])
            i -= 1; j -= 1
        elif dp[i-1][j] > dp[i][j-1]:
            i -= 1
        else:
            j -= 1
    return dp[n][m], ''.join(reversed(seq))

print(lcs("ABCBDAB","BDCAB"))  # (4, 'BCAB')
\end{lstlisting}

\textbf{Complexity.} $\bigO{nm}$ time and space (optimizable to $\bigO{m}$ space).

% ─────────────────────────────────────────────────────────────────────────────
\subsection{Interval DP}

\begin{algorithm}[Interval DP]
DP over contiguous intervals $[l, r]$, expanding from small to large. Iterate over interval length, then left endpoint, then split point $k$:
\[
  dp[l][r] = \min_{l \leq k < r} \bigl(dp[l][k] + dp[k{+}1][r] + \text{cost}(l,r)\bigr).
\]
\end{algorithm}

\textbf{Where to use.} Matrix chain multiplication; optimal BST; merging stones; bursting balloons; palindrome partitioning.

\textbf{Typical CP scenarios.} Merging adjacent elements with additive costs; splitting or grouping a sequence optimally.

\textbf{Example Problem.} Given $n$ piles of stones in a row, repeatedly merge two adjacent piles (cost = sum of the two). Minimize total cost.

\subsubsection{Python Implementation}
\begin{lstlisting}[caption={Interval DP -- merge stones and matrix chain}]
def merge_stones(stones):
    """Minimum cost to merge all adjacent piles into one."""
    n = len(stones)
    prefix = [0] * (n + 1)
    for i in range(n):
        prefix[i+1] = prefix[i] + stones[i]
    INF = float('inf')
    dp = [[0]*n for _ in range(n)]
    for length in range(2, n + 1):
        for l in range(n - length + 1):
            r = l + length - 1
            dp[l][r] = INF
            cost = prefix[r+1] - prefix[l]
            for k in range(l, r):
                dp[l][r] = min(dp[l][r],
                               dp[l][k] + dp[k+1][r] + cost)
    return dp[0][n-1]

def matrix_chain(dims):
    """dims = [d0,...,dn]. Matrix i is dims[i]*dims[i+1].
    Minimize scalar multiplications."""
    n = len(dims) - 1
    dp = [[0]*n for _ in range(n)]
    for length in range(2, n + 1):
        for i in range(n - length + 1):
            j = i + length - 1
            dp[i][j] = float('inf')
            for k in range(i, j):
                dp[i][j] = min(dp[i][j], dp[i][k] + dp[k+1][j]
                               + dims[i]*dims[k+1]*dims[j+1])
    return dp[0][n-1]

print(merge_stones([3,4,5,2]))  # 33
print(matrix_chain([10,30,5,60]))  # 4500
\end{lstlisting}

\textbf{Complexity.} $\bigO{n^3}$ time; $\bigO{n^2}$ space.

\begin{cpnote}
For $n > 500$, check if Knuth's optimization applies (quadrangle inequality) to reduce to $\bigO{n^2}$.
\end{cpnote}

% ─────────────────────────────────────────────────────────────────────────────
\subsection{Bitmask DP}

\subsubsection{Python Implementation}
\begin{lstlisting}[caption={TSP via Bitmask DP}]
def tsp(dist, n):
    """Minimum Hamiltonian cycle. dist[i][j] = cost."""
    INF = float('inf')
    dp = [[INF] * n for _ in range(1 << n)]
    dp[1][0] = 0  # Start at node 0, visited = {0}
    for mask in range(1 << n):
        for u in range(n):
            if not (mask >> u & 1): continue
            if dp[mask][u] == INF: continue
            for v in range(n):
                if mask >> v & 1: continue
                nmask = mask | (1 << v)
                dp[nmask][v] = min(dp[nmask][v], dp[mask][u] + dist[u][v])
    full = (1 << n) - 1
    ans = min(dp[full][u] + dist[u][0] for u in range(n))
    return ans
\end{lstlisting}

\textbf{Complexity.} $\bigO{2^n \cdot n^2}$.

% ─────────────────────────────────────────────────────────────────────────────
\subsection{Digit DP}

\subsubsection{Python Implementation}
\begin{lstlisting}[caption={Digit DP -- count numbers with digit sum divisible by k}]
from functools import lru_cache

def digit_dp(N, K):
    """Count integers in [1, N] with digit sum divisible by K."""
    digits = list(map(int, str(N)))
    n = len(digits)

    @lru_cache(maxsize=None)
    def dp(pos, rem, tight, started):
        if pos == n:
            return 1 if (rem == 0 and started) else 0
        limit = digits[pos] if tight else 9
        res = 0
        for d in range(0, limit + 1):
            new_rem = (rem + d) % K if (started or d > 0) else 0
            new_started = started or (d > 0)
            res += dp(pos+1, new_rem, tight and d==limit, new_started)
        return res

    return dp(0, 0, True, False)

print(digit_dp(100, 3))  # count in [1,100] with digit sum % 3 == 0
\end{lstlisting}

\textbf{Complexity.} $\bigO{\text{digits} \times K \times 10}$.

% ─────────────────────────────────────────────────────────────────────────────
\subsection{Tree DP}

\subsubsection{Python Implementation}
\begin{lstlisting}[caption={Tree DP -- maximum independent set on tree}]
def tree_max_independent_set(tree, n, root=0):
    """dp[v][0] = max set not including v, dp[v][1] = including v."""
    dp = [[0, 1] for _ in range(n)]  # [exclude, include]
    parent = [-1] * n
    order = []
    stack = [root]
    visited = [False] * n
    visited[root] = True
    while stack:
        u = stack.pop()
        order.append(u)
        for v in tree[u]:
            if not visited[v]:
                visited[v] = True
                parent[v] = u
                stack.append(v)
    for u in reversed(order):
        for v in tree[u]:
            if parent[v] == u:
                dp[u][0] += max(dp[v][0], dp[v][1])
                dp[u][1] += dp[v][0]
    return max(dp[root][0], dp[root][1])
\end{lstlisting}

% ─────────────────────────────────────────────────────────────────────────────
\subsection{DP with Prefix Sums}

\textbf{Pattern.} When $dp[i] = \sum_{j<i} f(dp[j])$, precompute prefix sums of $dp$ to answer range queries in $\bigO{1}$.

\subsubsection{Python Implementation}
\begin{lstlisting}[caption={DP with prefix sums -- count sequences}]
def count_sequences(n, mod=10**9+7):
    """Example: dp[i] = sum of dp[j] for all valid j < i."""
    dp = [0] * (n + 1)
    prefix = [0] * (n + 2)
    dp[0] = 1
    prefix[1] = 1
    for i in range(1, n + 1):
        # Example: valid j in [i//2, i-1]
        lo, hi = i // 2, i - 1
        dp[i] = (prefix[hi + 1] - prefix[lo]) % mod
        prefix[i + 1] = (prefix[i] + dp[i]) % mod
    return dp[n]
\end{lstlisting}

% ─────────────────────────────────────────────────────────────────────────────
\subsection{Meet-in-the-Middle DP}

\textbf{Idea.} Split the problem into two halves of size $n/2$. Solve each independently ($\bigO{2^{n/2}}$ each), then combine. Reduces $\bigO{2^n}$ to $\bigO{2^{n/2}}$.

\subsubsection{Python Implementation}
\begin{lstlisting}[caption={Meet-in-the-Middle -- subset sum}]
def meet_in_middle(arr, target):
    """Check if any subset sums to target."""
    n = len(arr)
    half = n // 2
    left, right = arr[:half], arr[half:]

    def all_subset_sums(a):
        sums = set()
        for mask in range(1 << len(a)):
            s = sum(a[i] for i in range(len(a)) if mask >> i & 1)
            sums.add(s)
        return sums

    left_sums = all_subset_sums(left)
    right_sums = all_subset_sums(right)
    for s in left_sums:
        if target - s in right_sums:
            return True
    return False
\end{lstlisting}

\textbf{Complexity.} $\bigO{2^{n/2} \log n}$.

% ─────────────────────────────────────────────────────────────────────────────
\subsection{State Compression DP}

\textbf{Concept.} Represent sets/subsets as bitmasks (integers). Enumerate all $2^n$ subsets efficiently.

\begin{lstlisting}[caption={Subset enumeration -- iterate all subsets of a mask}]
# Iterate all subsets of mask (including mask itself and 0)
def subsets(mask):
    s = mask
    while s > 0:
        yield s
        s = (s - 1) & mask
    yield 0

# Sum over subsets (SOS DP)
def sos_dp(f, n):
    """f[mask] = sum over subsets after SOS DP."""
    for i in range(n):
        for mask in range(1 << n):
            if mask >> i & 1:
                f[mask] += f[mask ^ (1 << i)]
    return f
\end{lstlisting}

% ─────────────────────────────────────────────────────────────────────────────
\subsection{DP Optimizations}

\begin{itemize}[noitemsep]
  \item \textbf{Divide and Conquer DP:} When $opt(i,j)$ is monotone. Reduces $\bigO{n^2}$ to $\bigO{n \log n}$.
  \item \textbf{Knuth's Optimization:} For interval DP when quadrangle inequality holds. Reduces $\bigO{n^3}$ to $\bigO{n^2}$.
  \item \textbf{Convex Hull Trick:} When $dp[i] = \min_j (f(j) + x_i \cdot g(j))$ and $x_i$ is monotone.
\end{itemize}

\subsubsection{Python Implementation}
\begin{lstlisting}[caption={Convex Hull Trick (Li Chao tree variant)}]
class LiChaoTree:
    def __init__(self, xs):
        self.n = len(xs)
        self.xs = xs
        self.tree = [None] * (4 * self.n)

    def add_line(self, m, b, node=1, lo=0, hi=None):
        if hi is None: hi = self.n - 1
        if lo > hi: return
        mid = (lo + hi) // 2
        x = self.xs[mid]
        new_line = (m, b)
        if self.tree[node] is None:
            self.tree[node] = new_line
            return
        cur = self.tree[node]
        if m * x + b < cur[0] * x + cur[1]:
            self.tree[node] = new_line
            new_line = cur
        m2, b2 = new_line
        if m2 * self.xs[lo] + b2 < self.tree[node][0]*self.xs[lo]+self.tree[node][1]:
            self.add_line(m2, b2, 2*node, lo, mid-1)
        else:
            self.add_line(m2, b2, 2*node+1, mid+1, hi)

    def query(self, x_idx, node=1, lo=0, hi=None):
        if hi is None: hi = self.n - 1
        if lo > hi: return float('inf')
        mid = (lo + hi) // 2
        best = float('inf')
        if self.tree[node]:
            m, b = self.tree[node]
            best = m * self.xs[x_idx] + b
        if x_idx <= mid:
            return min(best, self.query(x_idx, 2*node, lo, mid-1))
        else:
            return min(best, self.query(x_idx, 2*node+1, mid+1, hi))
\end{lstlisting}

\clearpage
% =============================================================================
\section{Greedy Algorithms}
% =============================================================================

\subsection{Interval Scheduling / Activity Selection}

\begin{algorithm}[Activity Selection]
Sort intervals by finish time. Greedily select each interval that starts after the last selected one finishes.
\end{algorithm}

\textbf{Why correct?} Exchange argument: any solution can be transformed to the greedy one without reducing the count.

\subsubsection{Python Implementation}
\begin{lstlisting}[caption={Activity Selection -- maximum non-overlapping intervals}]
def activity_selection(intervals):
    """intervals = [(start, end)]. Returns max non-overlapping count."""
    intervals.sort(key=lambda x: x[1])
    count = 0
    last_end = float('-inf')
    selected = []
    for s, e in intervals:
        if s >= last_end:
            count += 1
            last_end = e
            selected.append((s, e))
    return count, selected

print(activity_selection([(1,3),(2,5),(3,9),(6,8)]))  # 3
\end{lstlisting}

\textbf{Complexity.} $\bigO{n \log n}$.

% ─────────────────────────────────────────────────────────────────────────────
\subsection{Huffman Coding}

\begin{algorithm}[Huffman]
Greedy prefix-free encoding. Repeatedly merge two nodes of lowest frequency. Produces optimal variable-length codes.
\end{algorithm}

\subsubsection{Python Implementation}
\begin{lstlisting}[caption={Huffman Coding}]
import heapq

def huffman(freq):
    """freq = {char: frequency}. Returns code lengths."""
    heap = [[f, [c, ""]] for c, f in freq.items()]
    heapq.heapify(heap)
    while len(heap) > 1:
        lo = heapq.heappop(heap)
        hi = heapq.heappop(heap)
        for pair in lo[1:]: pair[1] = '0' + pair[1]
        for pair in hi[1:]: pair[1] = '1' + pair[1]
        heapq.heappush(heap, [lo[0]+hi[0]] + lo[1:] + hi[1:])
    codes = {c: code for c, code in sorted(heap[0][1:], key=lambda x: x[0])}
    return codes

freq = {'a':5,'b':9,'c':12,'d':13,'e':16,'f':45}
print(huffman(freq))
\end{lstlisting}

\textbf{Complexity.} $\bigO{n \log n}$.

% ─────────────────────────────────────────────────────────────────────────────
\subsection{Greedy Exchange Argument}

\textbf{Template.} To prove a greedy is optimal: assume an optimal solution differs from greedy; show swapping to match greedy doesn't worsen the solution. Repeat until solutions match.

\textbf{Example.} Job scheduling to minimize total weighted completion time: sort by $w_i / t_i$ descending.

\subsubsection{Python Implementation}
\begin{lstlisting}[caption={Weighted job scheduling -- minimize weighted completion time}]
def weighted_job_schedule(jobs):
    """jobs = [(time, weight)]. Minimize sum of weight * completion_time."""
    # Sort by weight/time descending (or time/weight ascending)
    jobs.sort(key=lambda x: x[0] / x[1])  # ascending time/weight
    total_cost = 0
    cur_time = 0
    for t, w in jobs:
        cur_time += t
        total_cost += w * cur_time
    return total_cost
\end{lstlisting}

% ─────────────────────────────────────────────────────────────────────────────
\subsection{Fractional Knapsack}

\subsubsection{Python Implementation}
\begin{lstlisting}[caption={Fractional Knapsack}]
def fractional_knapsack(items, W):
    """items = [(weight, value)]. Returns max value (fractions allowed)."""
    items.sort(key=lambda x: x[1]/x[0], reverse=True)
    total = 0.0
    remaining = W
    for w, v in items:
        if remaining <= 0: break
        take = min(w, remaining)
        total += take * (v / w)
        remaining -= take
    return total

print(fractional_knapsack([(10,60),(20,100),(30,120)], 50))  # 240.0
\end{lstlisting}

\textbf{Complexity.} $\bigO{n \log n}$.

\clearpage
% =============================================================================
\section{Data Structures}
% =============================================================================

\subsection{Segment Tree}

\begin{algorithm}[Segment Tree]
Binary tree where each node stores aggregate (sum/min/max) of a range. Supports point update and range query in $\bigO{\log n}$.
\end{algorithm}

\subsubsection{Python Implementation}
\begin{lstlisting}[caption={Segment Tree -- range sum with point update}]
class SegTree:
    def __init__(self, n):
        self.n = n
        self.tree = [0] * (4 * n)

    def update(self, pos, val, node=1, lo=0, hi=None):
        if hi is None: hi = self.n - 1
        if lo == hi:
            self.tree[node] += val
            return
        mid = (lo + hi) // 2
        if pos <= mid: self.update(pos, val, 2*node, lo, mid)
        else:          self.update(pos, val, 2*node+1, mid+1, hi)
        self.tree[node] = self.tree[2*node] + self.tree[2*node+1]

    def query(self, l, r, node=1, lo=0, hi=None):
        if hi is None: hi = self.n - 1
        if r < lo or hi < l: return 0
        if l <= lo and hi <= r: return self.tree[node]
        mid = (lo + hi) // 2
        return (self.query(l, r, 2*node, lo, mid) +
                self.query(l, r, 2*node+1, mid+1, hi))

st = SegTree(8)
for i, v in enumerate([1,3,5,7,9,11,13,15]):
    st.update(i, v)
print(st.query(1, 5))  # 3+5+7+9+11 = 35
\end{lstlisting}

\textbf{Complexity.} $\bigO{\log n}$ update and query; $\bigO{n}$ build.

% ─────────────────────────────────────────────────────────────────────────────
\subsection{Lazy Propagation}

\subsubsection{Python Implementation}
\begin{lstlisting}[caption={Segment Tree with Lazy Propagation -- range add, range sum}]
class LazySegTree:
    def __init__(self, arr):
        self.n = len(arr)
        self.tree = [0] * (4 * self.n)
        self.lazy = [0] * (4 * self.n)
        self._build(arr, 1, 0, self.n - 1)

    def _build(self, arr, node, lo, hi):
        if lo == hi:
            self.tree[node] = arr[lo]; return
        mid = (lo + hi) // 2
        self._build(arr, 2*node, lo, mid)
        self._build(arr, 2*node+1, mid+1, hi)
        self.tree[node] = self.tree[2*node] + self.tree[2*node+1]

    def _push(self, node, lo, hi):
        if self.lazy[node]:
            mid = (lo + hi) // 2
            self.tree[2*node]   += self.lazy[node] * (mid - lo + 1)
            self.tree[2*node+1] += self.lazy[node] * (hi - mid)
            self.lazy[2*node]   += self.lazy[node]
            self.lazy[2*node+1] += self.lazy[node]
            self.lazy[node] = 0

    def update(self, l, r, val, node=1, lo=0, hi=None):
        if hi is None: hi = self.n - 1
        if r < lo or hi < l: return
        if l <= lo and hi <= r:
            self.tree[node] += val * (hi - lo + 1)
            self.lazy[node] += val; return
        self._push(node, lo, hi)
        mid = (lo + hi) // 2
        self.update(l, r, val, 2*node, lo, mid)
        self.update(l, r, val, 2*node+1, mid+1, hi)
        self.tree[node] = self.tree[2*node] + self.tree[2*node+1]

    def query(self, l, r, node=1, lo=0, hi=None):
        if hi is None: hi = self.n - 1
        if r < lo or hi < l: return 0
        if l <= lo and hi <= r: return self.tree[node]
        self._push(node, lo, hi)
        mid = (lo + hi) // 2
        return (self.query(l, r, 2*node, lo, mid) +
                self.query(l, r, 2*node+1, mid+1, hi))
\end{lstlisting}

% ─────────────────────────────────────────────────────────────────────────────
\subsection{Fenwick Tree (BIT)}

\begin{algorithm}[Binary Indexed Tree]
Efficiently computes prefix sums. Update and query in $\bigO{\log n}$ with very low constants.
\end{algorithm}

\subsubsection{Python Implementation}
\begin{lstlisting}[caption={Fenwick Tree -- prefix sum, range sum, 2D BIT}]
class BIT:
    def __init__(self, n):
        self.n = n
        self.tree = [0] * (n + 1)

    def update(self, i, delta=1):
        while i <= self.n:
            self.tree[i] += delta
            i += i & (-i)

    def query(self, i):
        s = 0
        while i > 0:
            s += self.tree[i]
            i -= i & (-i)
        return s

    def range_query(self, l, r):
        return self.query(r) - self.query(l - 1)

    def kth(self, k):
        """Find k-th element (1-indexed)."""
        pos = 0
        log = self.n.bit_length()
        for i in range(log, -1, -1):
            npos = pos + (1 << i)
            if npos <= self.n and self.tree[npos] < k:
                pos = npos
                k -= self.tree[npos]
        return pos + 1

# 2D BIT
class BIT2D:
    def __init__(self, n, m):
        self.n, self.m = n, m
        self.tree = [[0]*(m+1) for _ in range(n+1)]

    def update(self, x, y, delta=1):
        i = x
        while i <= self.n:
            j = y
            while j <= self.m:
                self.tree[i][j] += delta
                j += j & (-j)
            i += i & (-i)

    def query(self, x, y):
        s, i = 0, x
        while i > 0:
            j = y
            while j > 0:
                s += self.tree[i][j]
                j -= j & (-j)
            i -= i & (-i)
        return s
\end{lstlisting}

\textbf{Complexity.} $\bigO{\log n}$ per operation.

% ─────────────────────────────────────────────────────────────────────────────
\subsection{Sparse Table}

\begin{algorithm}[Sparse Table]
Precompute range-minimum (or idempotent aggregates) queries. $\bigO{n \log n}$ build, $\bigO{1}$ query.
\end{algorithm}

\subsubsection{Python Implementation}
\begin{lstlisting}[caption={Sparse Table for RMQ (range minimum query)}]
import math

class SparseTable:
    def __init__(self, arr):
        n = len(arr)
        LOG = max(1, math.floor(math.log2(n)) + 1) if n > 1 else 1
        self.table = [arr[:]]
        for k in range(1, LOG):
            prev = self.table[k-1]
            cur = [min(prev[i], prev[i + (1<<(k-1))])
                   if i + (1<<k) - 1 < n else prev[i]
                   for i in range(n)]
            self.table.append(cur)
        self.log2 = [0] * (n + 1)
        for i in range(2, n + 1):
            self.log2[i] = self.log2[i//2] + 1

    def query(self, l, r):
        k = self.log2[r - l + 1]
        return min(self.table[k][l], self.table[k][r - (1<<k) + 1])

st = SparseTable([2,4,3,1,6,7,8,9,1,7])
print(st.query(2, 7))  # min of [3,1,6,7,8,9] = 1
\end{lstlisting}

\textbf{Complexity.} $\bigO{n \log n}$ preprocess; $\bigO{1}$ query.

% ─────────────────────────────────────────────────────────────────────────────
\subsection{Trie}

\subsubsection{Python Implementation}
\begin{lstlisting}[caption={Trie -- insert, search, prefix count}]
class TrieNode:
    __slots__ = ['children', 'count', 'end']
    def __init__(self):
        self.children = {}
        self.count = 0
        self.end = 0

class Trie:
    def __init__(self):
        self.root = TrieNode()

    def insert(self, word):
        node = self.root
        for c in word:
            if c not in node.children:
                node.children[c] = TrieNode()
            node = node.children[c]
            node.count += 1
        node.end += 1

    def search(self, word):
        node = self.root
        for c in word:
            if c not in node.children: return False
            node = node.children[c]
        return node.end > 0

    def starts_with(self, prefix):
        node = self.root
        for c in prefix:
            if c not in node.children: return 0
            node = node.children[c]
        return node.count

# Binary Trie for XOR problems
class BinaryTrie:
    def __init__(self, bits=30):
        self.root = [None, None]
        self.BITS = bits

    def insert(self, num):
        node = self.root
        for i in range(self.BITS, -1, -1):
            bit = (num >> i) & 1
            if not node[bit]: node[bit] = [None, None]
            node = node[bit]

    def max_xor(self, num):
        node = self.root
        res = 0
        for i in range(self.BITS, -1, -1):
            bit = (num >> i) & 1
            want = 1 - bit
            if node[want]:
                res |= (1 << i)
                node = node[want]
            elif node[bit]:
                node = node[bit]
            else:
                break
        return res
\end{lstlisting}

\clearpage
% ─────────────────────────────────────────────────────────────────────────────
\subsection{Monotonic Stack}

\textbf{Use.} Next greater element; largest rectangle in histogram; stock span problem.

\subsubsection{Python Implementation}
\begin{lstlisting}[caption={Monotonic Stack -- next greater element and histogram}]
def next_greater(arr):
    n = len(arr)
    result = [-1] * n
    stack = []  # Stores indices
    for i in range(n):
        while stack and arr[stack[-1]] < arr[i]:
            result[stack.pop()] = arr[i]
        stack.append(i)
    return result

def largest_rectangle_histogram(heights):
    n = len(heights)
    stack = []
    max_area = 0
    for i in range(n + 1):
        h = heights[i] if i < n else 0
        while stack and heights[stack[-1]] > h:
            height = heights[stack.pop()]
            width = i if not stack else i - stack[-1] - 1
            max_area = max(max_area, height * width)
        stack.append(i)
    return max_area

print(largest_rectangle_histogram([2,1,5,6,2,3]))  # 10
\end{lstlisting}

% ─────────────────────────────────────────────────────────────────────────────
\subsection{Monotonic Queue (Deque)}

\textbf{Use.} Sliding window minimum/maximum in $\bigO{n}$.

\subsubsection{Python Implementation}
\begin{lstlisting}[caption={Monotonic Deque -- sliding window minimum}]
from collections import deque

def sliding_window_min(arr, k):
    dq = deque()  # Stores indices; front = min index
    result = []
    for i, x in enumerate(arr):
        while dq and arr[dq[-1]] >= x:
            dq.pop()
        dq.append(i)
        if dq[0] < i - k + 1:
            dq.popleft()
        if i >= k - 1:
            result.append(arr[dq[0]])
    return result

print(sliding_window_min([1,3,-1,-3,5,3,6,7], 3))
# [-1,-3,-3,-3,3,3]
\end{lstlisting}

% ─────────────────────────────────────────────────────────────────────────────
\subsection{Priority Queue (Heap)}

\subsubsection{Python Implementation}
\begin{lstlisting}[caption={Priority Queue patterns in Python}]
import heapq

# Min-heap (default)
heap = []
heapq.heappush(heap, 5)
heapq.heappush(heap, 1)
heapq.heappush(heap, 3)
print(heapq.heappop(heap))  # 1

# Max-heap: negate values
max_heap = []
heapq.heappush(max_heap, -5)
heapq.heappush(max_heap, -1)
print(-heapq.heappop(max_heap))  # 5

# K-th smallest in a stream
def kth_largest(stream, k):
    heap = []
    for x in stream:
        heapq.heappush(heap, x)
        if len(heap) > k:
            heapq.heappop(heap)
    return heap[0]

# Lazy deletion heap
class LazyHeap:
    def __init__(self):
        self.heap = []
        self.removed = {}
    def push(self, val):
        heapq.heappush(self.heap, val)
    def remove(self, val):
        self.removed[val] = self.removed.get(val, 0) + 1
    def pop(self):
        while self.heap and self.removed.get(self.heap[0], 0) > 0:
            v = heapq.heappop(self.heap)
            self.removed[v] -= 1
        return heapq.heappop(self.heap)
\end{lstlisting}

% ─────────────────────────────────────────────────────────────────────────────
\subsection{Hash Map / Hashing}

\subsubsection{Python Implementation}
\begin{lstlisting}[caption={Hashing patterns for competitive programming}]
from collections import defaultdict, Counter

# Frequency counting
def most_frequent(arr):
    return Counter(arr).most_common(1)[0]

# Group by hash
def group_anagrams(words):
    groups = defaultdict(list)
    for w in words:
        groups[tuple(sorted(w))].append(w)
    return list(groups.values())

# Polynomial hashing (anti-hack: use double hash)
class PolyHash:
    def __init__(self, s, base=131, mod=10**9+7):
        n = len(s)
        self.mod = mod
        self.h = [0] * (n + 1)
        self.pw = [1] * (n + 1)
        for i, c in enumerate(s):
            self.h[i+1] = (self.h[i] * base + ord(c)) % mod
            self.pw[i+1] = self.pw[i] * base % mod

    def get(self, l, r):  # 0-indexed, inclusive
        return (self.h[r+1] - self.h[l] * self.pw[r-l+1]) % self.mod
\end{lstlisting}

\clearpage
% =============================================================================
\section{String Algorithms}
% =============================================================================

\subsection{KMP Algorithm}

\begin{algorithm}[Knuth-Morris-Pratt]
Find all occurrences of pattern $P$ in text $T$ in $\bigO{|T|+|P|}$ using failure function.
\end{algorithm}

\subsubsection{Python Implementation}
\begin{lstlisting}[caption={KMP -- failure function and pattern matching}]
def kmp_fail(pattern):
    n = len(pattern)
    fail = [0] * n
    j = 0
    for i in range(1, n):
        while j > 0 and pattern[i] != pattern[j]:
            j = fail[j-1]
        if pattern[i] == pattern[j]:
            j += 1
        fail[i] = j
    return fail

def kmp_search(text, pattern):
    """Returns list of start indices of pattern in text."""
    fail = kmp_fail(pattern)
    result = []
    j = 0
    for i, c in enumerate(text):
        while j > 0 and c != pattern[j]:
            j = fail[j-1]
        if c == pattern[j]:
            j += 1
        if j == len(pattern):
            result.append(i - len(pattern) + 1)
            j = fail[j-1]
    return result

print(kmp_search("ababcababc", "ababc"))  # [0, 5]
\end{lstlisting}

\textbf{Complexity.} $\bigO{|T|+|P|}$.

% ─────────────────────────────────────────────────────────────────────────────
\subsection{Z-Algorithm}

\subsubsection{Python Implementation}
\begin{lstlisting}[caption={Z-Function -- z[i] = longest match of s[i:] with prefix}]
def z_function(s):
    n = len(s)
    z = [0] * n
    z[0] = n
    l = r = 0
    for i in range(1, n):
        if i < r:
            z[i] = min(r - i, z[i - l])
        while i + z[i] < n and s[z[i]] == s[i + z[i]]:
            z[i] += 1
        if i + z[i] > r:
            l, r = i, i + z[i]
    return z

def z_search(text, pattern):
    """Find pattern in text using Z-algorithm."""
    s = pattern + '$' + text
    z = z_function(s)
    m = len(pattern)
    return [i - m - 1 for i in range(m+1, len(s)) if z[i] == m]

print(z_search("ababcababc", "ababc"))  # [0, 5]
\end{lstlisting}

\textbf{Complexity.} $\bigO{n}$.

% ─────────────────────────────────────────────────────────────────────────────
\subsection{Suffix Array}

\subsubsection{Python Implementation}
\begin{lstlisting}[caption={Suffix Array -- $O(n \log^2 n)$ construction}]
def build_suffix_array(s):
    """O(n log^2 n) rank-pair sorting construction."""
    n = len(s)
    sa = list(range(n))
    rank = [ord(c) for c in s]
    tmp = [0] * n
    k = 1
    while k < n:
        key = lambda i: (rank[i], rank[i+k] if i+k < n else -1)
        sa.sort(key=key)
        tmp[sa[0]] = 0
        for i in range(1, n):
            tmp[sa[i]] = tmp[sa[i-1]]
            if key(sa[i]) != key(sa[i-1]):
                tmp[sa[i]] += 1
        rank = tmp[:]
        if rank[sa[-1]] == n - 1: break
        k <<= 1
    return sa

def build_lcp(s, sa):
    """Kasai's algorithm: LCP array in O(n)."""
    n = len(s)
    rank = [0] * n
    for i, v in enumerate(sa):
        rank[v] = i
    lcp = [0] * n
    h = 0
    for i in range(n):
        if rank[i] > 0:
            j = sa[rank[i]-1]
            while i+h < n and j+h < n and s[i+h] == s[j+h]:
                h += 1
            lcp[rank[i]] = h
            if h > 0: h -= 1
    return lcp

s = "banana"
sa = build_suffix_array(s)
lcp = build_lcp(s, sa)
print(sa)   # [5,3,1,0,4,2]
print(lcp)  # [0,1,3,0,0,2]
\end{lstlisting}

\textbf{Complexity.} Build: $\bigO{n \log^2 n}$; LCP: $\bigO{n}$.

% ─────────────────────────────────────────────────────────────────────────────
\subsection{Rolling Hash / Rabin-Karp}

\subsubsection{Python Implementation}
\begin{lstlisting}[caption={Rabin-Karp rolling hash for pattern matching}]
def rabin_karp(text, pattern, base=131, mod=10**9+7):
    n, m = len(text), len(pattern)
    if m > n: return []
    pm = bm = 0
    power = 1
    for i in range(m):
        pm = (pm * base + ord(pattern[i])) % mod
        bm = (bm * base + ord(text[i])) % mod
        if i > 0: power = power * base % mod
    result = []
    if pm == bm and text[:m] == pattern:
        result.append(0)
    for i in range(1, n - m + 1):
        bm = (bm - ord(text[i-1]) * power % mod + mod) % mod
        bm = (bm * base + ord(text[i+m-1])) % mod
        if bm == pm and text[i:i+m] == pattern:
            result.append(i)
    return result
\end{lstlisting}

\textbf{Complexity.} $\bigO{nm}$ worst case; $\bigO{n+m}$ expected.

% ─────────────────────────────────────────────────────────────────────────────
\subsection{Manacher's Algorithm}

\begin{algorithm}[Manacher]
Finds longest palindromic substring (and all palindrome radii) in $\bigO{n}$.
\end{algorithm}

\subsubsection{Python Implementation}
\begin{lstlisting}[caption={Manacher's Algorithm}]
def manacher(s):
    """Returns p[] where p[i] = radius of palindrome centered at i in transformed string."""
    t = '#' + '#'.join(s) + '#'
    n = len(t)
    p = [0] * n
    c = r = 0
    for i in range(n):
        mirror = 2*c - i
        if i < r:
            p[i] = min(r - i, p[mirror])
        while i + p[i] + 1 < n and i - p[i] - 1 >= 0 and t[i+p[i]+1] == t[i-p[i]-1]:
            p[i] += 1
        if i + p[i] > r:
            c, r = i, i + p[i]
    return p

def longest_palindrome(s):
    p = manacher(s)
    max_len, center = max((v, i) for i, v in enumerate(p))
    start = (center - max_len) // 2
    return s[start:start + max_len]

print(longest_palindrome("babad"))  # "bab" or "aba"
\end{lstlisting}

\textbf{Complexity.} $\bigO{n}$.

% ─────────────────────────────────────────────────────────────────────────────
\subsection{Trie-based Matching}

Already shown in Section 6 (Trie). For Aho-Corasick (multi-pattern matching):

\subsubsection{Python Implementation}
\begin{lstlisting}[caption={Aho-Corasick multi-pattern matching}]
from collections import deque

def aho_corasick(patterns, text):
    # Build trie
    goto = [{}]
    fail = [0]
    output = [[]]

    for pat in patterns:
        cur = 0
        for c in pat:
            if c not in goto[cur]:
                goto[cur][c] = len(goto)
                goto.append({})
                fail.append(0)
                output.append([])
            cur = goto[cur][c]
        output[cur].append(pat)

    # Build failure links via BFS
    q = deque()
    for c, s in goto[0].items():
        fail[s] = 0
        q.append(s)
    while q:
        r = q.popleft()
        for c, s in goto[r].items():
            q.append(s)
            state = fail[r]
            while state and c not in goto[state]:
                state = fail[state]
            fail[s] = goto[state].get(c, 0)
            if fail[s] == s: fail[s] = 0
            output[s] += output[fail[s]]

    # Search
    cur = 0
    results = []
    for i, c in enumerate(text):
        while cur and c not in goto[cur]:
            cur = fail[cur]
        cur = goto[cur].get(c, 0)
        for pat in output[cur]:
            results.append((i - len(pat) + 1, pat))
    return results
\end{lstlisting}

\clearpage
% =============================================================================
\section{Geometry}
% =============================================================================

\subsection{Cross Product \& Dot Product}

\begin{definition}
For 2D vectors $\vec{a}=(a_x,a_y)$ and $\vec{b}=(b_x,b_y)$:
\[
  \vec{a} \times \vec{b} = a_x b_y - a_y b_x \quad \text{(scalar, sign = orientation)}
\]
\[
  \vec{a} \cdot \vec{b} = a_x b_x + a_y b_y \quad \text{(scalar, sign = angle relationship)}
\]
\end{definition}

\subsubsection{Python Implementation}
\begin{lstlisting}[caption={Cross product, dot product, orientation test}]
def cross(O, A, B):
    """Signed area of triangle OAB * 2. >0: CCW, <0: CW, =0: collinear."""
    return (A[0]-O[0])*(B[1]-O[1]) - (A[1]-O[1])*(B[0]-O[0])

def dot(O, A, B):
    return (A[0]-O[0])*(B[0]-O[0]) + (A[1]-O[1])*(B[1]-O[1])

def orientation(P, Q, R):
    val = cross(P, Q, R)
    if val == 0: return 0   # Collinear
    return 1 if val > 0 else -1  # 1=CCW, -1=CW

def dist2(A, B):
    """Squared distance (avoids float)."""
    return (A[0]-B[0])**2 + (A[1]-B[1])**2
\end{lstlisting}

% ─────────────────────────────────────────────────────────────────────────────
\subsection{Convex Hull (Andrew's Monotone Chain)}

\subsubsection{Python Implementation}
\begin{lstlisting}[caption={Convex Hull -- Andrew's Monotone Chain O(n log n)}]
def convex_hull(points):
    points = sorted(set(points))
    if len(points) <= 1: return points
    lower = []
    for p in points:
        while len(lower) >= 2 and cross(lower[-2], lower[-1], p) <= 0:
            lower.pop()
        lower.append(p)
    upper = []
    for p in reversed(points):
        while len(upper) >= 2 and cross(upper[-2], upper[-1], p) <= 0:
            upper.pop()
        upper.append(p)
    return lower[:-1] + upper[:-1]

def cross(O, A, B):
    return (A[0]-O[0])*(B[1]-O[1]) - (A[1]-O[1])*(B[0]-O[0])

pts = [(0,0),(1,1),(2,2),(0,2),(2,0)]
hull = convex_hull(pts)
print(hull)  # [(0,0),(2,0),(2,2),(0,2)]
\end{lstlisting}

\textbf{Complexity.} $\bigO{n \log n}$.

% ─────────────────────────────────────────────────────────────────────────────
\subsection{Shoelace Formula \& Polygon Area}

\[
  \text{Area} = \frac{1}{2}\left|\sum_{i=0}^{n-1}(x_i y_{i+1} - x_{i+1} y_i)\right|
\]

\subsubsection{Python Implementation}
\begin{lstlisting}[caption={Shoelace formula for polygon area}]
def polygon_area(pts):
    """Signed area * 2 (integer). Positive = CCW."""
    n = len(pts)
    area2 = 0
    for i in range(n):
        x1, y1 = pts[i]
        x2, y2 = pts[(i+1) % n]
        area2 += x1 * y2 - x2 * y1
    return abs(area2) / 2

print(polygon_area([(0,0),(4,0),(4,3),(0,3)]))  # 12.0
\end{lstlisting}

% ─────────────────────────────────────────────────────────────────────────────
\subsection{Pick's Theorem}

\begin{theorem}[Pick's Theorem]
For a lattice polygon with integer vertices:
\[
  A = I + \frac{B}{2} - 1,
\]
where $A$ = area, $I$ = interior lattice points, $B$ = boundary lattice points.
\end{theorem}

\subsubsection{Python Implementation}
\begin{lstlisting}[caption={Pick's Theorem -- interior lattice points}]
from math import gcd

def boundary_points(pts):
    """Count lattice points on boundary of polygon."""
    n = len(pts)
    b = 0
    for i in range(n):
        x1, y1 = pts[i]
        x2, y2 = pts[(i+1) % n]
        b += gcd(abs(x2-x1), abs(y2-y1))
    return b

def interior_points(pts):
    area2 = int(sum((pts[i][0]*pts[(i+1)%len(pts)][1] -
                     pts[(i+1)%len(pts)][0]*pts[i][1])
                    for i in range(len(pts))))
    area2 = abs(area2)
    B = boundary_points(pts)
    return (area2 - B + 2) // 2  # Pick: I = A - B/2 + 1
\end{lstlisting}

% ─────────────────────────────────────────────────────────────────────────────
\subsection{Line Intersection}

\subsubsection{Python Implementation}
\begin{lstlisting}[caption={Line segment intersection}]
def on_segment(P, Q, R):
    """Is R on segment PQ (assuming collinear)?"""
    return (min(P[0],Q[0]) <= R[0] <= max(P[0],Q[0]) and
            min(P[1],Q[1]) <= R[1] <= max(P[1],Q[1]))

def segments_intersect(P1, Q1, P2, Q2):
    d1 = cross(P2, Q2, P1)
    d2 = cross(P2, Q2, Q1)
    d3 = cross(P1, Q1, P2)
    d4 = cross(P1, Q1, Q2)
    if ((d1>0 and d2<0 or d1<0 and d2>0) and
        (d3>0 and d4<0 or d3<0 and d4>0)):
        return True
    if d1==0 and on_segment(P2, Q2, P1): return True
    if d2==0 and on_segment(P2, Q2, Q1): return True
    if d3==0 and on_segment(P1, Q1, P2): return True
    if d4==0 and on_segment(P1, Q1, Q2): return True
    return False

def cross(O, A, B):
    return (A[0]-O[0])*(B[1]-O[1]) - (A[1]-O[1])*(B[0]-O[0])

def line_intersection(P1, P2, P3, P4):
    """Intersection point of lines P1P2 and P3P4 (if exists)."""
    d1x, d1y = P2[0]-P1[0], P2[1]-P1[1]
    d2x, d2y = P4[0]-P3[0], P4[1]-P3[1]
    denom = d1x*d2y - d1y*d2x
    if denom == 0: return None  # Parallel
    t = ((P3[0]-P1[0])*d2y - (P3[1]-P1[1])*d2x) / denom
    return (P1[0] + t*d1x, P1[1] + t*d1y)
\end{lstlisting}

% ─────────────────────────────────────────────────────────────────────────────
\subsection{Point in Polygon}

\subsubsection{Python Implementation}
\begin{lstlisting}[caption={Point-in-polygon ray casting}]
def point_in_polygon(point, polygon):
    """Ray casting algorithm. Returns True if point is inside polygon."""
    x, y = point
    n = len(polygon)
    inside = False
    px, py = polygon[-1]
    for qx, qy in polygon:
        if ((py > y) != (qy > y) and
                x < (qx - px) * (y - py) / (qy - py) + px):
            inside = not inside
        px, py = qx, qy
    return inside

# Winding number (more robust)
def winding_number(point, polygon):
    x, y = point
    wn = 0
    n = len(polygon)
    for i in range(n):
        x1, y1 = polygon[i]
        x2, y2 = polygon[(i+1) % n]
        if y1 <= y:
            if y2 > y:
                if cross((x1,y1),(x2,y2),(x,y)) > 0:
                    wn += 1
        else:
            if y2 <= y:
                if cross((x1,y1),(x2,y2),(x,y)) < 0:
                    wn -= 1
    return wn  # 0 = outside

def cross(O, A, B):
    return (A[0]-O[0])*(B[1]-O[1]) - (A[1]-O[1])*(B[0]-O[0])
\end{lstlisting}

% ─────────────────────────────────────────────────────────────────────────────
\subsection{Distance Formulas}

\textbf{Mathematical formulations.}
\begin{itemize}[noitemsep]
  \item \textbf{Euclidean:} $d = \sqrt{(x_2-x_1)^2 + (y_2-y_1)^2}$
  \item \textbf{Manhattan:} $d = |x_2-x_1| + |y_2-y_1|$ (grid, taxicab)
  \item \textbf{Chebyshev:} $d = \max(|x_2-x_1|, |y_2-y_1|)$ (king moves)
  \item \textbf{Chebyshev $\leftrightarrow$ Manhattan trick:} $(x,y) \to (x+y, x-y)$ converts Chebyshev to Manhattan and vice versa.
\end{itemize}

\begin{cputility}
Grid problems; closest pair; geometry distance queries; metric space arguments.
\end{cputility}

\subsubsection{Python Implementation}
\begin{lstlisting}[caption={Distance formulas and closest pair}]
import math

def euclidean(P, Q):
    return math.hypot(P[0]-Q[0], P[1]-Q[1])

def manhattan(P, Q):
    return abs(P[0]-Q[0]) + abs(P[1]-Q[1])

def chebyshev(P, Q):
    return max(abs(P[0]-Q[0]), abs(P[1]-Q[1]))

def chebyshev_to_manhattan(points):
    """Transform so Chebyshev dist becomes Manhattan."""
    return [(x+y, x-y) for x, y in points]

def closest_pair(pts):
    """Closest pair of points in O(n log n)."""
    pts.sort()
    def rec(P):
        if len(P) <= 3:
            return min((euclidean(P[i],P[j]), P[i], P[j])
                       for i in range(len(P))
                       for j in range(i+1, len(P)))
        mid = len(P) // 2
        mx = P[mid][0]
        d, *pair = min(rec(P[:mid]), rec(P[mid:]))
        strip = [p for p in P if abs(p[0]-mx) < d]
        strip.sort(key=lambda p: p[1])
        for i in range(len(strip)):
            j = i + 1
            while j < len(strip) and strip[j][1]-strip[i][1] < d:
                nd = euclidean(strip[i], strip[j])
                if nd < d:
                    d = nd; pair = [strip[i], strip[j]]
                j += 1
        return (d, *pair)
    return rec(pts)

print(closest_pair([(0,0),(1,1),(3,4),(2,2)]))  # closest pair
\end{lstlisting}

\textbf{Complexity.} Individual distances: $\bigO{1}$; Closest pair: $\bigO{n \log n}$.

\clearpage
% =============================================================================
\section{Probability \& Mathematics}
% =============================================================================

\subsection{Expected Value \& Linearity of Expectation}

\begin{theorem}[Linearity of Expectation]
For any random variables $X_1, \ldots, X_n$ (not necessarily independent):
\[
  \mathbb{E}\left[\sum_{i=1}^n X_i\right] = \sum_{i=1}^n \mathbb{E}[X_i].
\]
\end{theorem}

\textbf{Intuition.} Decompose complex random variables into simple indicator variables.

\textbf{Classic application.} Expected number of comparisons in randomized algorithms; expected cycle length in random permutations.

\subsubsection{Python Implementation}
\begin{lstlisting}[caption={Expected value via linearity -- hat-check problem}]
def expected_hats_returned(n):
    """Expected fixed points in a random permutation of n elements."""
    # By linearity: E[X] = sum E[X_i] = n * (1/n) = 1
    return 1.0  # Always 1, regardless of n!

def expected_steps(P, start, target):
    """Expected steps from start to absorbing target state.
    P[i] = [(next_state, probability), ...]. Gaussian elim."""
    n = len(P)
    # System: E[i] - sum_j p_ij*E[j] = 1 (i!=target); E[target]=0
    A = [[0.0]*(n+1) for _ in range(n)]
    for i in range(n):
        if i == target:
            A[i][i] = 1.0
        else:
            A[i][i] = 1.0; A[i][n] = 1.0
            for j, p in P[i]:
                A[i][j] -= p
    for col in range(n):  # Forward elimination
        for row in range(col + 1, n):
            if abs(A[col][col]) < 1e-12: continue
            f = A[row][col] / A[col][col]
            for k in range(n + 1):
                A[row][k] -= f * A[col][k]
    E = [0.0] * n         # Back substitution
    for i in range(n - 1, -1, -1):
        E[i] = A[i][n]
        for j in range(i + 1, n):
            E[i] -= A[i][j] * E[j]
        if abs(A[i][i]) > 1e-12:
            E[i] /= A[i][i]
    return E[start]

# Gambler's ruin: states 0..4, absorb at 0 and 4
# P[i] = [(i-1, 0.5), (i+1, 0.5)] for 1<=i<=3
# expected_steps(P, 2, 4)  # expected steps from 2 to 4
\end{lstlisting}

% ─────────────────────────────────────────────────────────────────────────────
\subsection{Bayes' Theorem}

\begin{theorem}[Bayes' Theorem]
For hypotheses $H_1, \ldots, H_n$ forming a partition:
\[
  P(A \mid B) = \frac{P(B \mid A) \cdot P(A)}{P(B)}, \qquad
  P(H_i \mid B) = \frac{P(B \mid H_i)\, P(H_i)}{\sum_j P(B \mid H_j)\, P(H_j)}.
\]
\end{theorem}

\textbf{Intuition.} Bayes inverts conditional probabilities: given $P(\text{evidence} \mid \text{hypothesis})$, compute $P(\text{hypothesis} \mid \text{evidence})$.

\begin{cputility}
Probability problems with conditional information; interactive problems with Bayesian updates; classification scenarios.
\end{cputility}

\textbf{Example Problem.} Two bags: Bag~1 has 3 red, 2 blue balls; Bag~2 has 1 red, 4 blue. A bag is chosen uniformly, then a ball is drawn. Given the ball is red, what is the probability it came from Bag~1?

\subsubsection{Python Implementation}
\begin{lstlisting}[caption={Bayes' Theorem -- posterior computation}]
def bayes(prior, likelihood):
    """prior[i] = P(H_i), likelihood[i] = P(evidence | H_i).
    Returns posterior[i] = P(H_i | evidence)."""
    joint = [prior[i] * likelihood[i] for i in range(len(prior))]
    total = sum(joint)
    return [j / total for j in joint]

# Bag 1: 3R/5, Bag 2: 1R/5. Prior = [0.5, 0.5]
prior = [0.5, 0.5]
likelihood = [3/5, 1/5]
print(bayes(prior, likelihood))  # [0.75, 0.25]
\end{lstlisting}

\textbf{Complexity.} $\bigO{n}$ for $n$ hypotheses.

% ─────────────────────────────────────────────────────────────────────────────
\subsection{Randomized Algorithms}

\begin{definition}
A \textbf{Monte Carlo} algorithm gives the correct answer with high probability (may be wrong). A \textbf{Las Vegas} algorithm always gives the correct answer but has random running time.
\end{definition}

\subsubsection{Python Implementation}
\begin{lstlisting}[caption={Randomized algorithms -- Miller-Rabin primality and random hash}]
import random

def miller_rabin(n, k=20):
    """Probabilistic primality test. Error prob < 4^(-k)."""
    if n < 2: return False
    if n == 2 or n == 3: return True
    if n % 2 == 0: return False
    r, d = 0, n - 1
    while d % 2 == 0:
        r += 1; d //= 2
    for _ in range(k):
        a = random.randrange(2, n - 1)
        x = pow(a, d, n)
        if x == 1 or x == n - 1: continue
        for _ in range(r - 1):
            x = x * x % n
            if x == n - 1: break
        else:
            return False
    return True

def deterministic_miller_rabin(n):
    """Deterministic for n < 3.2*10^18."""
    witnesses = [2,3,5,7,11,13,17,19,23,29,31,37]
    if n < 2: return False
    if n in witnesses: return True
    if any(n % w == 0 for w in witnesses): return False
    r, d = 0, n - 1
    while d % 2 == 0:
        r += 1; d //= 2
    for a in witnesses:
        if a >= n: continue
        x = pow(a, d, n)
        if x == 1 or x == n - 1: continue
        for _ in range(r - 1):
            x = x * x % n
            if x == n - 1: break
        else:
            return False
    return True

print(deterministic_miller_rabin(998244353))  # True (prime)
\end{lstlisting}

\textbf{Complexity.} $\bigO{k \log^2 n}$ per test.

% ─────────────────────────────────────────────────────────────────────────────
\subsection{Randomized Techniques in CP}

\begin{itemize}[noitemsep]
  \item \textbf{Randomized hashing:} Use random base/mod to avoid adversarial anti-hash tests.
  \item \textbf{Random shuffling:} Shuffle input before applying deterministic algorithm (Quicksort).
  \item \textbf{Random sampling:} Estimate quantities via Monte Carlo.
\end{itemize}

\subsubsection{Python Implementation}
\begin{lstlisting}[caption={Anti-hack double hashing for strings}]
import random

class AntiHackHash:
    def __init__(self, s):
        MOD1 = (1 << 61) - 1  # Mersenne prime
        MOD2 = 10**9 + 9
        BASE1 = random.randint(131, 1000000)
        BASE2 = random.randint(137, 1000000)
        n = len(s)
        self.h1 = [0]*(n+1); self.h2 = [0]*(n+1)
        self.p1 = [1]*(n+1); self.p2 = [1]*(n+1)
        for i, c in enumerate(s):
            self.h1[i+1] = (self.h1[i]*BASE1 + ord(c)) % MOD1
            self.h2[i+1] = (self.h2[i]*BASE2 + ord(c)) % MOD2
            self.p1[i+1] = self.p1[i]*BASE1 % MOD1
            self.p2[i+1] = self.p2[i]*BASE2 % MOD2
        self.M1, self.M2 = MOD1, MOD2

    def get(self, l, r):  # 0-indexed inclusive
        v1 = (self.h1[r+1] - self.h1[l]*self.p1[r-l+1]) % self.M1
        v2 = (self.h2[r+1] - self.h2[l]*self.p2[r-l+1]) % self.M2
        return (v1, v2)
\end{lstlisting}

\clearpage
% =============================================================================
\section{CP Patterns}
% =============================================================================

\subsection{Sliding Window}

\textbf{Pattern.} Maintain a window $[l, r]$ that expands right and contracts left. $\bigO{n}$ for fixed-size or variable-size windows.

\subsubsection{Python Implementation}
\begin{lstlisting}[caption={Sliding Window -- maximum subarray sum of length k}]
def max_subarray_fixed(arr, k):
    n = len(arr)
    win_sum = sum(arr[:k])
    best = win_sum
    for i in range(k, n):
        win_sum += arr[i] - arr[i-k]
        best = max(best, win_sum)
    return best

def min_subarray_sum_geq_s(arr, s):
    """Minimum length subarray with sum >= s."""
    n = len(arr)
    best = float('inf')
    cur = l = 0
    for r in range(n):
        cur += arr[r]
        while cur >= s:
            best = min(best, r - l + 1)
            cur -= arr[l]
            l += 1
    return best if best != float('inf') else 0

def longest_no_repeat(s):
    """Longest substring without repeating characters."""
    last = {}
    l = best = 0
    for r, c in enumerate(s):
        if c in last and last[c] >= l:
            l = last[c] + 1
        best = max(best, r - l + 1)
        last[c] = r
    return best
\end{lstlisting}

% ─────────────────────────────────────────────────────────────────────────────
\subsection{Two Pointers}

\subsubsection{Python Implementation}
\begin{lstlisting}[caption={Two Pointers -- pair sum, three sum}]
def two_sum_sorted(arr, target):
    """Find pair summing to target in sorted array."""
    l, r = 0, len(arr)-1
    while l < r:
        s = arr[l] + arr[r]
        if s == target: return (l, r)
        elif s < target: l += 1
        else: r -= 1
    return None

def three_sum(arr):
    """All unique triplets summing to 0."""
    arr.sort()
    n = len(arr)
    result = []
    for i in range(n-2):
        if i > 0 and arr[i] == arr[i-1]: continue
        l, r = i+1, n-1
        while l < r:
            s = arr[i] + arr[l] + arr[r]
            if s == 0:
                result.append([arr[i], arr[l], arr[r]])
                while l < r and arr[l]==arr[l+1]: l+=1
                while l < r and arr[r]==arr[r-1]: r-=1
                l+=1; r-=1
            elif s < 0: l+=1
            else: r-=1
    return result
\end{lstlisting}

% ─────────────────────────────────────────────────────────────────────────────
\subsection{Binary Search on Answer}

\textbf{Pattern.} If the answer satisfies a monotone predicate (false, false, \ldots, true, true), binary search on the answer space.

\subsubsection{Python Implementation}
\begin{lstlisting}[caption={Binary search on answer -- minimize maximum}]
def binary_search_answer(lo, hi, check):
    """Find smallest x in [lo, hi] where check(x) is True."""
    while lo < hi:
        mid = (lo + hi) // 2
        if check(mid):
            hi = mid
        else:
            lo = mid + 1
    return lo

# Example: koko eats bananas
def min_eating_speed(piles, H):
    import math
    def can_finish(k):
        return sum(math.ceil(p/k) for p in piles) <= H
    return binary_search_answer(1, max(piles), can_finish)

# Example: allocate minimum pages
def allocate_books(pages, m):
    def feasible(max_pages):
        students = 1; cur = 0
        for p in pages:
            if p > max_pages: return False
            if cur + p > max_pages:
                students += 1; cur = 0
            cur += p
        return students <= m
    return binary_search_answer(max(pages), sum(pages), feasible)
\end{lstlisting}

% ─────────────────────────────────────────────────────────────────────────────
\subsection{Prefix Sums}

\subsubsection{Python Implementation}
\begin{lstlisting}[caption={Prefix Sums -- 1D, 2D, and difference arrays}]
def prefix_sum_1d(arr):
    n = len(arr)
    pre = [0] * (n + 1)
    for i, v in enumerate(arr):
        pre[i+1] = pre[i] + v
    return pre

def range_sum(pre, l, r):  # 0-indexed inclusive
    return pre[r+1] - pre[l]

# 2D prefix sum
def prefix_sum_2d(grid):
    n, m = len(grid), len(grid[0])
    pre = [[0]*(m+1) for _ in range(n+1)]
    for i in range(n):
        for j in range(m):
            pre[i+1][j+1] = (grid[i][j] + pre[i][j+1] +
                              pre[i+1][j] - pre[i][j])
    return pre

def rect_sum(pre, r1, c1, r2, c2):
    return (pre[r2+1][c2+1] - pre[r1][c2+1] -
            pre[r2+1][c1] + pre[r1][c1])
\end{lstlisting}

% ─────────────────────────────────────────────────────────────────────────────
\subsection{Difference Arrays}

\subsubsection{Python Implementation}
\begin{lstlisting}[caption={Difference Arrays -- range update in O(1)}]
def range_add(diff, l, r, val):
    diff[l] += val
    if r + 1 < len(diff): diff[r+1] -= val

def apply_diff(diff):
    arr = diff[:]
    for i in range(1, len(arr)):
        arr[i] += arr[i-1]
    return arr

# Example: add val to [l,r] for multiple queries, then retrieve
n = 10
diff = [0] * (n + 1)
range_add(diff, 1, 4, 3)
range_add(diff, 2, 6, -1)
result = apply_diff(diff)
print(result[:n])  # [0,3,2,2,2,-1,-1,0,0,0]
\end{lstlisting}

% ─────────────────────────────────────────────────────────────────────────────
\subsection{Coordinate Compression}

\subsubsection{Python Implementation}
\begin{lstlisting}[caption={Coordinate Compression}]
def compress(arr):
    """Map values to indices 0..k-1."""
    sorted_unique = sorted(set(arr))
    rank = {v: i for i, v in enumerate(sorted_unique)}
    return [rank[v] for v in arr], sorted_unique

def compress_2d(points):
    xs = sorted(set(p[0] for p in points))
    ys = sorted(set(p[1] for p in points))
    rx = {v: i for i, v in enumerate(xs)}
    ry = {v: i for i, v in enumerate(ys)}
    return [(rx[x], ry[y]) for x, y in points]

arr = [100, 200, 150, 300, 200]
comp, vals = compress(arr)
print(comp)  # [0, 2, 1, 3, 2]
print(vals)  # [100, 150, 200, 300]
\end{lstlisting}

% ─────────────────────────────────────────────────────────────────────────────
\subsection{Mo's Algorithm}

\begin{algorithm}[Mo's Algorithm]
Answer offline range queries in $\bigO{(n+q)\sqrt{n}}$ by sorting queries by block.
\end{algorithm}

\subsubsection{Python Implementation}
\begin{lstlisting}[caption={Mo's Algorithm for offline range queries}]
def mo_algorithm(arr, queries):
    """queries = [(l, r, idx)]. Returns answers[]."""
    n = len(arr)
    block = max(1, int(n**0.5))
    queries = sorted(queries,
                     key=lambda q: (q[0]//block,
                                    q[1] if (q[0]//block)%2==0 else -q[1]))
    cnt = {}
    cur_ans = 0
    ans = [0] * len(queries)

    def add(x):
        nonlocal cur_ans
        cnt[x] = cnt.get(x, 0) + 1
        # Update cur_ans (problem-specific)

    def remove(x):
        nonlocal cur_ans
        cnt[x] -= 1
        if cnt[x] == 0: del cnt[x]
        # Update cur_ans

    l = r = 0
    add(arr[0])
    for ql, qr, qi in queries:
        while r < qr: r += 1; add(arr[r])
        while l > ql: l -= 1; add(arr[l])
        while r > qr: remove(arr[r]); r -= 1
        while l < ql: remove(arr[l]); l += 1
        ans[qi] = cur_ans
    return ans
\end{lstlisting}

\textbf{Complexity.} $\bigO{(n+q)\sqrt{n}}$.

% ─────────────────────────────────────────────────────────────────────────────
\subsection{Offline Queries}

\textbf{Pattern.} Sort queries and process them in a convenient order (by right endpoint, by value, etc.) using an online data structure (BIT, segment tree).

\subsubsection{Python Implementation}
\begin{lstlisting}[caption={Offline range distinct count using BIT}]
def offline_distinct(arr, queries):
    """Count distinct elements in [l,r] for each query."""
    n = len(arr)
    # Process queries offline sorted by right endpoint
    qs = sorted(enumerate(queries), key=lambda x: x[1][1])
    bit = [0] * (n + 1)

    def update(i, v):
        while i <= n:
            bit[i] += v; i += i & -i

    def query(i):
        s = 0
        while i > 0:
            s += bit[i]; i -= i & -i
        return s

    last = {}
    ans = [0] * len(queries)
    qi = 0
    for r in range(n):
        x = arr[r]
        if x in last:
            update(last[x]+1, -1)
        update(r+1, 1)
        last[x] = r
        while qi < len(qs) and qs[qi][1][1] == r:
            orig_i, (l, _) = qs[qi]
            ans[orig_i] = query(r+1) - query(l)
            qi += 1
    return ans
\end{lstlisting}

% ─────────────────────────────────────────────────────────────────────────────
\subsection{Bitmasking}

\subsubsection{Python Implementation}
\begin{lstlisting}[caption={Bitmasking tricks for competitive programming}]
# Common bitmask operations
def bits(mask, n):
    """Iterate set bits."""
    return [i for i in range(n) if mask >> i & 1]

def lowest_bit(x): return x & (-x)

def pop_lowest(x): return x & (x - 1)

def submasks(mask):
    s = mask
    while s:
        yield s
        s = (s - 1) & mask

# Check if exactly k bits set
def has_k_bits(x, k): return bin(x).count('1') == k

# Turn on/off bits
def set_bit(x, i):    return x | (1 << i)
def clear_bit(x, i):  return x & ~(1 << i)
def toggle_bit(x, i): return x ^ (1 << i)
def check_bit(x, i):  return bool(x >> i & 1)

# Iterate all subsets of size k
from itertools import combinations
def subsets_of_size(n, k):
    for bits in combinations(range(n), k):
        yield sum(1 << b for b in bits)
\end{lstlisting}

% ─────────────────────────────────────────────────────────────────────────────
\subsection{Subset Iteration Trick (SOS DP)}

\subsubsection{Python Implementation}
\begin{lstlisting}[caption={SOS (Sum over Subsets) DP}]
def sos_dp(f, n):
    """f[mask] = sum of f[submask] for all submasks of mask.
    O(n * 2^n)."""
    for i in range(n):
        for mask in range(1 << n):
            if mask >> i & 1:
                f[mask] += f[mask ^ (1 << i)]
    return f

# Example: f[mask] = number of submasks with property P
n = 4
f = [1, 0, 1, 0, 1, 0, 0, 0, 0, 0, 1, 0, 0, 0, 0, 0]  # 2^4
sos_dp(f, n)
# f[mask] now = count of submasks that have property P
\end{lstlisting}

\textbf{Complexity.} $\bigO{n \cdot 2^n}$.

% ─────────────────────────────────────────────────────────────────────────────
\subsection{Meet-in-the-Middle}

\textbf{Pattern.} Split the input in half ($n/2$ each), enumerate all possibilities for each half (e.g., $2^{n/2}$ subsets), then combine efficiently using sorting + binary search or two pointers.

\begin{cputility}
Reduces exponential brute force from $\bigO{2^n}$ to $\bigO{2^{n/2} \log n}$. Critical when $n \leq 40$.
\end{cputility}

\textbf{Where to use.} Subset sum with $n \leq 40$; 4-SUM; knapsack with large $n$ but small halves.

\subsubsection{Python Implementation}
\begin{lstlisting}[caption={Meet-in-the-Middle -- count subsets summing to target}]
from bisect import bisect_left, bisect_right

def mitm_count(arr, target):
    """Count subsets of arr that sum to target. n <= 40."""
    n = len(arr)
    half = n // 2
    left, right = arr[:half], arr[half:]

    def all_sums(a):
        sums = []
        for mask in range(1 << len(a)):
            sums.append(sum(a[i] for i in range(len(a)) if mask >> i & 1))
        return sorted(sums)

    left_s = all_sums(left)
    right_s = all_sums(right)
    count = 0
    for s in left_s:
        need = target - s
        count += bisect_right(right_s, need) - bisect_left(right_s, need)
    return count

print(mitm_count([1,2,3,4,5,6], 10))  # subsets summing to 10
\end{lstlisting}

\textbf{Complexity.} $\bigO{2^{n/2} \log n}$.

% ─────────────────────────────────────────────────────────────────────────────
\subsection{Fast I/O Template}

\textbf{Why it matters.} Python's default \texttt{input()} is slow. For problems with $n > 10^5$, always use fast I/O to avoid TLE.

\subsubsection{Python Implementation}
\begin{lstlisting}[caption={Fast I/O template for competitive programming}]
import sys
input = sys.stdin.readline

def main():
    # Read entire input at once (fastest)
    data = sys.stdin.buffer.read().split()
    idx = 0
    def rd(): nonlocal idx; idx += 1; return int(data[idx-1])

    n = rd()
    arr = [rd() for _ in range(n)]

    # Process...
    ans = sum(arr)

    # Fast output
    sys.stdout.write(str(ans) + '\n')

main()

# Alternative: for multiple test cases
def solve():
    import sys
    input = sys.stdin.readline
    out = []
    T = int(input())
    for _ in range(T):
        n = int(input())
        arr = list(map(int, input().split()))
        out.append(str(sum(arr)))
    print('\n'.join(out))
\end{lstlisting}

\begin{cpnote}
In Python, \texttt{sys.stdin.buffer.read()} into a byte buffer and splitting once is $5{-}10\times$ faster than repeated \texttt{input()} calls. Always use this for $n > 10^5$.
\end{cpnote}

\vspace{2cm}

\begin{mdframed}[backgroundcolor=theorembg, linecolor=accent, linewidth=2pt,
  innerleftmargin=20pt, innerrightmargin=20pt, innertopmargin=16pt, innerbottommargin=16pt]
\begin{center}
{\Large\bfseries\color{accent} Quick Reference Complexity Table}
\end{center}
\vspace{0.5em}
\begin{tabular}{>{\bfseries}p{4.5cm} p{3.5cm} p{4cm}}
\toprule
Algorithm & Time & Space \\
\midrule
Sieve of Eratosthenes & $\bigO{n \log \log n}$ & $\bigO{n}$ \\
GCD (Euclidean) & $\bigO{\log n}$ & $\bigO{1}$ \\
Binary Exponentiation & $\bigO{\log n}$ & $\bigO{1}$ \\
BFS / DFS & $\bigO{V+E}$ & $\bigO{V}$ \\
Dijkstra & $\bigO{(V+E)\log V}$ & $\bigO{V+E}$ \\
Bellman-Ford & $\bigO{VE}$ & $\bigO{V}$ \\
Floyd-Warshall & $\bigO{V^3}$ & $\bigO{V^2}$ \\
Kruskal & $\bigO{E \log E}$ & $\bigO{V}$ \\
Tarjan/Kosaraju SCC & $\bigO{V+E}$ & $\bigO{V}$ \\
Dinic Max Flow & $\bigO{V^2 E}$ & $\bigO{V+E}$ \\
LCA Binary Lifting & $\bigO{N \log N}$ / $\bigO{\log N}$ & $\bigO{N \log N}$ \\
Segment Tree & $\bigO{n}$ build, $\bigO{\log n}$ q/u & $\bigO{n}$ \\
Fenwick Tree & $\bigO{\log n}$ q/u & $\bigO{n}$ \\
Sparse Table & $\bigO{n \log n}$ / $\bigO{1}$ & $\bigO{n \log n}$ \\
KMP & $\bigO{n+m}$ & $\bigO{m}$ \\
Z-Algorithm & $\bigO{n}$ & $\bigO{n}$ \\
Suffix Array & $\bigO{n \log^2 n}$ & $\bigO{n}$ \\
Manacher & $\bigO{n}$ & $\bigO{n}$ \\
Convex Hull & $\bigO{n \log n}$ & $\bigO{n}$ \\
Closest Pair & $\bigO{n \log n}$ & $\bigO{n}$ \\
Interval DP & $\bigO{n^3}$ & $\bigO{n^2}$ \\
0/1 Knapsack & $\bigO{nW}$ & $\bigO{W}$ \\
LIS & $\bigO{n \log n}$ & $\bigO{n}$ \\
Bitmask DP (TSP) & $\bigO{2^n n^2}$ & $\bigO{2^n n}$ \\
Mo's Algorithm & $\bigO{(n+q)\sqrt{n}}$ & $\bigO{n+q}$ \\
SOS DP & $\bigO{n \cdot 2^n}$ & $\bigO{2^n}$ \\
Miller-Rabin & $\bigO{k \log^2 n}$ & $\bigO{1}$ \\
\bottomrule
\end{tabular}
\end{mdframed}

\end{document}
